\section{Overview of
\emph{Ricci Flow with Surgery on Three-Manifolds} \cite{Perelman2}} \label{overview2}

The paper studies Ricci flow with surgery on compact oriented $3$-manifolds.
Sections II.1--II.3 concern ancient solutions, the standard solution, and the
first singular time; Sections II.4--II.5 define surgery and establish its
existence; Sections II.6--II.8 analyze the long-time behavior.

\subsection{Ancient solutions, standard solutions, and first singularities}
\dcref{II.1--II.3}

Every three-dimensional $\kappa$-solution has an asymptotic soliton, obtained
as a rescaled limit as $t\to-\infty$. Such a soliton is a shrinking round
cylinder $\R\times S^2$, its $\Z_2$-quotient
$\R\times_{\Z_2}S^2$, or a finite quotient of the shrinking round $S^3$.
Every compact $\kappa$-solution is a finite quotient of the round shrinking
$S^3$, or is diffeomorphic to $S^3$ or $\R P^3$. There is a universal
$\kappa_0>0$ such that every $\kappa$-solution is either a finite quotient of
the round shrinking $S^3$ or a $\kappa_0$-solution; consequently its scalar
curvature has universal derivative bounds.

The standard solution is a Ricci flow on $\R^3$ whose initial metric is a
capped-off half-cylinder. It exists on the maximal interval $[0,1)$ and its
curvature becomes singular everywhere as $t\to1$.

At a first singular time $T$, set
\[
\Omega=\{x\in M:\limsup_{t\to T^-}|\Rm(x,t)|<\infty\}.
\]
Then $\Omega$ is open and
\[
x\in M-\Omega\quad\Longleftrightarrow\quad
\lim_{t\to T^-}R(x,t)=\infty.
\]
If $\Omega=\emptyset$, then $M$ is diffeomorphic to
$S^1\times S^2$, $\R P^3\#\R P^3$, or a finite isometric quotient of
$S^3$. Just before $T$, the flow then consists entirely of
high-scalar-curvature regions. If $\Omega\ne\emptyset$, it has a limit metric $\overline g$ with
scalar curvature $\overline R$. For sufficiently small $\rho>0$, define
\[
\Omega_\rho=\{x\in\Omega:\overline R(x)\le\rho^{-2}\}.
\]
This is compact. A component of $\Omega$ disjoint from $\Omega_\rho$ is a
capped $\epsilon$-horn or a double $\epsilon$-horn; a component meeting
$\Omega_\rho$ has finitely many ends, each an $\epsilon$-horn. Surgery
truncates every such horn and caps it by a $3$-ball, discards components
disjoint from $\Omega_\rho$, and produces a manifold $M'$. The original $M$
is recovered from the connected sum of the components of $M'$ together with
finitely many $S^1\times S^2$ and $\R P^3$ factors.

\subsection{Surgery construction and existence}
\dcref{II.4--II.5}

A Ricci flow with surgery is a sequence of smooth $3$-dimensional Ricci flows
on adjacent time intervals such that adjacent final and initial slices share a
compact $3$-dimensional submanifold with boundary. It satisfies Hamilton--Ivey
pinching and the canonical neighborhood assumption: if
$R(x,t)\ge r(t)^{-2}$, then after rescaling the neighborhood of $(x,t)$ is
$\epsilon$-close to a neighborhood in a $\kappa$-solution or in a time slice
of the standard solution. Here $r(t)$ is decreasing and $\epsilon$ is a small
universal constant.

The parameters are a decreasing canonical-neighborhood scale $r(t)$, a
nonincreasing function $\delta(t)\to0$, the truncation scale
\[
\rho(t)=\delta(t)r(t),
\]
and the surgery scale $h$. If a singular time is $T$, then for every horn in
a component meeting $\Omega_{\rho(T)}$, the horn self-improves: for every
$\delta>0$ there is $h<\delta\rho(T)$ such that
$\overline R(x)\ge h^{-2}$ implies that $x$ is the center of a
$\delta$-neck. Every horn point is the center of an $\epsilon$-neck, which is
rescaled $\epsilon$-close to
$[-\epsilon^{-1},\epsilon^{-1}]\times S^2$. Cutting at the center of such a neck and gluing a rescaled
truncated standard solution defines the surgery; components disjoint from
$\Omega_{\rho(T)}$ are discarded. A volume-change estimate gives finitely
many surgeries in every finite time interval, though not necessarily finitely
many surgeries altogether.

The pinching condition holds before the first surgery by Hamilton--Ivey and
can be preserved by the surgery. For a normalized initial metric there is a
universal $C$ such that the flow is smooth with bounded curvature on $[0,C]$;
choosing $r(t)$ sufficiently small there makes the canonical neighborhood
condition vacuous. On the intervals
\[
[2^{j-1}\epsilon,2^j\epsilon],\qquad j\ge1,
\]
choose decreasing sequences $r_j$, $\kappa_j$, and
$\overline\delta_j$ so that
\[
\delta(t)\le\overline\delta_j\quad\text{on }[2^{j-1}\epsilon,2^j\epsilon],
\qquad
r(t)=r_j\quad\text{on }[2^{j-1}\epsilon,2^j\epsilon).
\]
Then the flow is $\kappa_j$-noncollapsed at scales less than $\epsilon$ and
satisfies the canonical neighborhood assumption on that interval. The
parameters are chosen inductively; at the $(i+1)$-st step only
$\overline\delta_i$ from the preceding interval is also redefined.

For each $r>0$, there is a $\kappa>0$ and
$\overline\delta=\overline\delta(r)>0$ such that if
$T\in[2^i\epsilon,2^{i+1}\epsilon]$, the proposition holds through
$2^i\epsilon$, the canonical neighborhood assumption with scale $r$ holds
on $[2^i\epsilon,T)$, and $\delta(t)\le\overline\delta$ on
$[2^{i-1}\epsilon,T)$, then the flow is $\kappa$-noncollapsed at scales less
than $\epsilon$. The proof uses ${\mathcal L}$-lengths: for sufficiently
small $\overline\delta$, competitive minimizing curves avoid the long thin
surgery necks, so the argument of I.7 applies on $[2^i\epsilon,T)$.

The induction proposition follows by contradiction. If it failed, take
$r^\al\to0$, $\overline\delta^\al\to0$, and a first counterexample
$(x^\al,t^\al)$ in $[2^i\epsilon,2^{i+1}\epsilon]$. The preceding lemma gives
noncollapsing up to this point. If a suitable backward parabolic neighborhood
contains no surgery, rescaling yields a limiting $\kappa$-solution as in
Theorem I.12.1; if surgeries interfere, the standard-solution persistence
estimate of Lemma II.4.5 puts $(x^\al,t^\al)$ in a canonical neighborhood.
Both alternatives contradict the choice of the counterexample. For a
simply-connected initial manifold, finite extinction follows from
\cite{Colding-Minicozzi,Colding-Minicozzi2,Perelman3}, yielding the
Poincare Conjecture.

\subsection{Long-time behavior}
\dcref{II.6--II.8}

Sections II.6 and II.8 establish back-and-forth curvature estimates analogous
to Theorems I.12.2 and I.12.3, with surgery terms included. For small $w$ and
large $t$, the $w$-thick part $M^+(w,t)$, with
\[
\widehat g(t)=t^{-1}g(t),
\]
approaches the $w$-thick part of a finite collection
$\{(H_i,x_i)\}_{i=1}^k$ of pointed finite-volume $3$-manifolds of constant
sectional curvature $-\frac14$. The cuspidal tori are incompressible. The
minimal compressing-disk argument uses an area differential inequality that
would force the area negative for large $t$.

The $w$-thin part $M^-(w,t)$ is a graph manifold when $w$ is small and $t$ is
large. The geometric hypothesis is that for each $x$ there is a radius
$\rho(x)$ with
\[
\vol B(x,\rho)\le w\rho^3,
\qquad
\sec\ge-\rho^{-2}\quad\text{on }B(x,\rho).
\]

Let $\lambda_1(g)=\lambda_1(-4\triangle+R)$. The manifold admits a metric
with $\lambda_1(g)>0$ iff it admits a metric with positive scalar curvature,
iff it is diffeomorphic to a connected sum of $S^1\times S^2$'s and round
quotients of $S^3$. Otherwise set
\[
\overline\lambda=\sup_g\lambda_1(g)\vol(M,g)^{2/3}.
\]
If $\overline\lambda=0$, then $M$ is a graph manifold; if
$\overline\lambda<0$, its geometric decomposition contains a nonempty
hyperbolic piece of total volume
\[
\left(-\frac23\overline\lambda\right)^{3/2}.
\]
These conclusions use monotonicity of
$\lambda(g(t))\vol(M,g(t))^{2/3}$ when it is nonpositive. The surgery
parameters can be chosen so this quantity is arbitrarily close to
nondecreasing.

\section{Notation and terminology}
\dcref{II}
\label{notationterminology}

$B(x,t,r)$ denotes the open metric ball of radius $r$ at time $t$, centered at
$x$. The parabolic neighborhood $P(x,t,r,\Delta t)$ consists of points
$(x',t')$ with $x'\in B(x,t,r)$ and $t'$ between $t$ and $t+\Delta t$.

\begin{definition}[smooth closeness and epsilon-necks]
  \label{def:kl-notation-terminology-smooth-closeness-epsilon-necks}
  \source{kleiner-lott}
  \dcref{Definition 58.1}
  \group{Kleiner--Lott Surgery Foundations}
 \label{closenessdef}
Two Riemannian manifolds $(M_1,g_1)$ and $(M_2,g_2)$ have distance at most
$\epsilon$ in the $C^N$ topology if a diffeomorphism
$\phi:M_2\to M_1$ satisfies
\[
\sum_{|I|\le N}\frac{1}{|I|!}
\|\nabla^I(\phi^*g_1-g_2)\|_\infty\le\epsilon.
\]
An open set in a Riemannian $3$-manifold is an $\eps$-neck if, after
rescaling, it is at distance less than $\eps$ in the
$C^{[1/\epsilon]+1}$ topology from
$S^2\times I$, where $S^2$ is round with scalar curvature $1$ and
$|I|>2\eps^{-1}$. The pointed version uses the center of $I$ as basepoint.
In spacetime, $C^N$-closeness includes time derivatives. A strong $\eps$-neck
is, after parabolic rescaling and time translation, $\eps$-close on
$[-1,0]$ to the product Ricci flow whose final slice is the same round
cylinder.
\end{definition}

Containment in an $\epsilon$-neck is open in the pointed
$C^{[1/\epsilon]+1}$ topology. If
$f:S^2\times I\to U$ is an $\epsilon$-approximation, a cross-sectional sphere
is $f(S^2\times\{\lambda\})$ for
$\lambda\in(-\epsilon^{-1},\epsilon^{-1})$. A curve crossing both ends of
the neck meets every cross-sectional sphere; a minimizing geodesic does so
exactly once when $\epsilon$ is sufficiently small. The strong-neck
parabolic neighborhood is backward, namely
$P(x,t,\epsilon^{-1}r,-r^2)$, and its rescaled time interval is $[-1,0]$.

In the next definition, $\I$ is an open interval and $B^3$ is an open ball.

\begin{definition}[epsilon-tubes horns and caps]
  \label{def:kl-notation-terminology-epsilon-tubes-horns-caps}
  \source{kleiner-lott}
  \dcref{Definition 2}
  \group{Kleiner--Lott Surgery Foundations}

A metric on $S^2\times\I$ whose points all lie in $\epsilon$-necks is an
${\em\epsilon}$-tube, an ${\em\epsilon}$-horn, or a double
${\em\epsilon}$-horn according as scalar curvature is bounded at both ends,
bounded at one end and tends to infinity at the other, or tends to infinity at
both ends. A metric on $B^3$ or $\R P^3-\overline{B^3}$ whose complement of
a compact subset is contained in $\epsilon$-necks is an $\epsilon$-cap or a
capped $\epsilon$-horn according as scalar curvature stays bounded or tends to
infinity on the end.
\end{definition}

An $\epsilon$-tube is exemplified by
$S^2\times(-\epsilon^{-1},\epsilon^{-1})$ with the product metric. A model
$\epsilon$-horn is the metric
\begin{equation}
g=dr^2+\frac{1}{8\ln(1/r)}r^2d\theta^2
\end{equation}
on $(0,R)\times S^2$, with $R=1$; this is the rotationally symmetric
neckpinch model of \cite{Angenent-Knopf}. Under
\[
s=\sqrt{8\ln(1/r_0)}\left(\frac r{r_0}-1\right),
\]
it satisfies
\begin{equation}
\frac{8\ln(1/r_0)}{r_0^2}g
=ds^2+\left(1+\frac{1}{\sqrt{8\ln(1/r_0)}}
\left(2+\frac{1}{\ln(1/r_0)}\right)s+O(s^2)\right)d\theta^2.
\end{equation}
For small $r_0$, taking
$\epsilon\sim(\ln(1/r_0))^{-1/4}$ makes the region
$|s|<\epsilon^{-1}$ $\epsilon$-biLipschitz close to the standard cylinder.
An $\epsilon$-cap is an $\epsilon$-tube capped by a $3$-ball or
$\R P^3-\overline{B^3}$; a capped $\epsilon$-horn is defined similarly.

\begin{remark}[positivity for notation and terminology]
  \label{rem:kl-notation-terminology-positivity-notation-terminology}
  \source{kleiner-lott}
  \dcref{Remark 58.5}
  \group{Kleiner--Lott Surgery Foundations}

\label{remepsilon}
Throughout, $\epsilon$ is a universal sufficiently small positive constant.
Statements involving $\epsilon$ are uniform in the other variables after
finitely many reductions of its size.
\end{remark}

\begin{lemma}[neck geometry lemma]
  \label{lem:kl-notation-terminology-neck-geometry-lemma}
  \source{kleiner-lott}
  \dcref{Lemma 4}
  \group{Kleiner--Lott Surgery Foundations}

If $U$ is an $\epsilon$-neck in an $\epsilon$-tube or horn and $S$ is a
cross-sectional sphere in $U$, then $S$ separates the two ends.
\end{lemma}
\begin{proof}
\sketch The lowest Ricci eigenspaces define a smooth line field transverse to
$S$. If $S$ did not separate the ends, it would bound an embedded $3$-ball,
contradicting extension of this line field over the ball.
\end{proof}

\section{Three-dimensional $\kappa$-solutions}
\dcref{II.1} \label{II.1}

We omit ``three-dimensional'' and ``oriented'' below. A $\kappa$-solution has
nonnegative sectional curvature. If curvature is not strictly positive, its
universal cover splits off a line and the solution is
$\R\times S^2$ or $\R\times_{\Z_2}S^2$. If curvature is strictly positive,
the compact case is a spherical space form \cite{Hamiltonnnnn}, while the
noncompact case is diffeomorphic to $\R^3$ \cite{Cheeger-Gromoll}. Let
$M_\epsilon$ denote the points that are not centers of $\epsilon$-necks.

\begin{lemma}[existence for kappa solutions]
  \label{lem:kl-three-dimensional-solutions-existence-kappa-solutions}
  \source{kleiner-lott}
  \dcref{Lemma 59.1}
  \group{Kleiner--Lott Surgery Foundations}
 \label{noncompactsoliton}
If a time slice of a noncompact $\kappa$-solution has
$M_\epsilon\ne\emptyset$, then there is a compact submanifold-with-boundary
$X\subset M$ containing $M_\epsilon$ such that
$X\cong\overline{B^3}$ or $\R P^3-B^3$, and
$M-\Int(X)\cong[0,\infty)\times S^2$.
\end{lemma}
\begin{proof}
\sketch In the nonpositive-curvature case the solution is
$\R\times_{\Z_2}S^2$. In the positive-curvature case, cover a ray from a
point of $M_\epsilon$ by overlapping necks, splice their projection maps into
a submersion onto a ray with fiber $S^2$, and choose a convex exhaustion level;
the compact side is a $3$-ball.
\end{proof}

\begin{lemma}[kappa solutions lemma]
  \label{lem:kl-three-dimensional-solutions-kappa-solutions-lemma}
  \source{kleiner-lott}
  \dcref{Lemma 59.2}
  \group{Kleiner--Lott Surgery Foundations}
 \label{onlyneck}
If a time slice of a $\kappa$-solution has $M_\epsilon=\emptyset$, then its
Ricci flow is the evolving round cylinder $\R\times S^2$.
\end{lemma}
\begin{proof}
\sketch The necks piece together to give a two-ended universal cover.
Toponogov splitting and the strong maximum principle split off an $\R$ factor;
the two-dimensional classification gives the evolving round cylinder, and
$\kappa$-noncollapsing excludes the quotient $S^1\times S^2$.
\end{proof}

A compact $\kappa$-solution with compact asymptotic soliton is a shrinking
quotient of the round $S^3$ \cite{Hamiltonnnnn}. The remaining compact case is
covered by the following lemma.

\begin{lemma}[compactness for kappa solutions]
  \label{lem:kl-three-dimensional-solutions-compactness-kappa-solutions}
  \source{kleiner-lott}
  \dcref{Lemma 59.3}
  \group{Kleiner--Lott Surgery Foundations}
 \label{S3RP3}
If a $\kappa$-solution $(M,g(\cdot))$ is compact and has a noncompact
asymptotic soliton, then $M$ is diffeomorphic to $S^3$ or $\R P^3$.
\end{lemma}
\begin{proof}
  \uses{cor:kl-necklike-behavior-at-infinity-three-dimensional-solution-strong-version-existence-kappa-solutions,lem:kl-three-dimensional-solutions-existence-kappa-solutions,thm:kl-compactness-space-three-dimensional-solutions-compactness-kappa-solutions}
\sketch The type-D alternative gives a universal upper bound for
$\max_M R\,\diam(M)^2$, ruling out a noncompact asymptotic soliton. In the
remaining case, rescale at separated points of $M_\epsilon(t_i)$ as
$t_i\to-\infty$. The pointed limit is a noncompact $\kappa$-solution, and the
corresponding compact regions are diffeomorphic to $\overline{B^3}$ or
$\R P^3-B^3$. The necks between them form an $S^2$-fibration, so $M$ is their
gluing along a sphere. Since $M$ has finite fundamental group, at most one
piece is projective, giving $S^3$ or $\R P^3$.
\end{proof}

Every ancient solution that is a $\kappa$-solution for some $\kappa$ is either
a $\kappa_0$-solution or a metric quotient of the round $S^3$.

\begin{lemma}[existence for kappa solutions]
  \label{lem:kl-three-dimensional-solutions-existence-kappa-solutions-4}
  \source{kleiner-lott}
  \dcref{Lemma 59.4}
  \group{Kleiner--Lott Surgery Foundations}
 \label{derivativeestimates}
There is a universal constant $\eta$ such that every ancient solution which is
a $\kappa$-solution for some $\kappa$ satisfies
\begin{equation} \label{derestimates}
|\nabla R|<\eta R^{3/2},\qquad |R_t|<\eta R^2.
\end{equation}
\end{lemma}
\begin{proof}
\sketch The round-sphere quotients satisfy the estimates directly. For
$\kappa_0$-solutions, rescale at the chosen point so that $R=1$ and apply the
compactness derivative bounds of Theorem \ref{thmI.11.7}.
\end{proof}

At points where $R\ne0$, this implies
\begin{equation}
\label{eqnderestimatesrinverse}
|\nabla(R^{-1/2})|<\frac\eta2,\qquad |(R^{-1})_t|<\eta.
\end{equation}

\begin{lemma}[existence for kappa solutions]
  \label{lem:kl-three-dimensional-solutions-existence-kappa-solutions-5}
  \source{kleiner-lott}
  \dcref{Lemma 59.7}
  \group{Kleiner--Lott Surgery Foundations}
 \label{nbhd}
For every sufficiently small $\epsilon>0$, there are
$C_1=C_1(\epsilon)$ and $C_2=C_2(\epsilon)$ such that for every point
$(x,t)$ in every $\kappa$-solution there is
\[
r\in[R(x,t)^{-1/2},C_1R(x,t)^{-1/2}]
\]
and a neighborhood $B$ with
$\overline{B(x,t,r)}\subset B\subset B(x,t,2r)$, satisfying one of:
\begin{itemize}
\item[(a)] $B$ is a strong $\epsilon$-neck;
\item[(b)] $B$ is an $\epsilon$-cap;
\item[(c)] $B$ is closed and diffeomorphic to $S^3$ or $\R P^3$;
\item[(d)] $B$ is a closed manifold of constant positive sectional curvature.
\end{itemize}
Moreover, the scalar curvature on $B$ at time $t$ lies between
$C_2^{-1}R(x,t)$ and $C_2R(x,t)$; in cases (a), (b), and (c),
\[
\vol(B)>C_2^{-1}R(x,t)^{-3/2};
\]
in case (b), an $\eps$-neck $U\subset B$ with compact complement lies at
distance at least $10000R(x,t)^{-1/2}$ from $x$; and in case (c), the
sectional curvature on $B$ is greater than $C_2^{-1}R(x,t)$.
\end{lemma}

\begin{proof}
  \uses{lem:kl-three-dimensional-solutions-existence-kappa-solutions,lem:kl-three-dimensional-solutions-kappa-solutions-lemma}
\sketch A spherical quotient is covered by taking a ball containing the whole
slice, for $r>\pi(R(x_i,t_i)/6)^{-1/2}$. Otherwise, a counterexample sequence
rescaled by $R(x_i,t_i)$ has a pointed $\kappa_0$-solution limit. A compact
limit satisfies the closed-manifold alternatives and their curvature and
volume bounds; a noncompact limit satisfies the neck/cap decomposition, with
the end-distance bound in the cap case. The subsidiary estimates pass to the
limit, a contradiction.
\end{proof}

\begin{remark}[structural remark on three-dimensional -solutions]
  \label{rem:kl-three-dimensional-solutions-structural-remark-three-dimensional-solutions}
  \source{kleiner-lott}
  \dcref{Remark 6}
  \group{Kleiner--Lott Surgery Foundations}

The lemma strengthens II.1.5 by requiring
$r\ge R(x,t)^{-1/2}$ rather than merely $r>0$.
\end{remark}

\section{Standard solutions}

\dcref{II.2} \label{II.2}

Fix a smooth rotationally symmetric complete metric $g_0$ on $\R^3$, obtained by
gluing a hemispherical-type cap to a half-infinite cylinder of scalar curvature
$1$. Assume that $g_0$ has nonnegative curvature operator and scalar curvature
at least $1$.

\begin{remark}[structural remark on standard solutions]
  \label{rem:kl-standard-solutions-structural-remark-standard-solutions}
  \source{kleiner-lott}
  \dcref{Remark 1}
  \group{Kleiner--Lott Surgery Foundations}

In Section \ref{surgery}, the initial metric $g_0$ is further specialized for
the surgery construction.
\end{remark}

\begin{definition}[standard solution definition]
  \label{def:kl-standard-solutions-standard-solution-definition}
  \source{kleiner-lott}
  \dcref{Definition 60.2}
  \group{Kleiner--Lott Surgery Foundations}

\label{defstandardsolution}
A Ricci flow $(\R^3,g(\cdot))$ on $[0,a)$ is a \emph{standard solution} if its
initial metric is $g_0$, $|\Rm|$ is bounded on every compact subinterval
$[0,a']\subset[0,a)$, and it has no extension with these properties to a
strictly longer time interval.
\end{definition}

\begin{lemma}[existence for kappa solutions]
  \label{lem:kl-standard-solutions-existence-kappa-solutions}
  \source{kleiner-lott}
  \uses{def:kl-standard-solutions-standard-solution-definition, thm:kl-canonical-neighborhood-theorem-compactness-canonical-neighborhoods}
  \dcref{Lemma 60.3}
  \group{Kleiner--Lott Surgery Foundations}

\label{lemstandard0}
If $(\R^3,g(\cdot))$ is a standard solution, then:

(1) The curvature operator is nonnegative.

(2) For small time, all curvature derivatives have bounds independent of the
standard solution.

(3) $\displaystyle\lim_{t\to a^-}\sup_{x\in\R^3}R(x,t)=\infty$.

(4) The flow is $\kappa$-noncollapsed at scales below $1$ on every time
interval contained in $[0,2]$, for a $\kappa$ depending only on $g_0$.

(5) For every $\Delta t>0$ there is $r_0>0$ such that, if $t_0\geq\Delta t$ and
$Q=R(x_0,t_0)\geq r_0^{-2}$, then the parabolic region
\[
\{(x,t):\dist_{t_0}^2(x,x_0)<(\epsilon Q)^{-1},\quad
t_0-(\epsilon Q)^{-1}\leq t\leq t_0\}
\]
is, after scaling by $Q$, $\epsilon$-close to the corresponding region of a
$\kappa$-solution.

Any flow satisfying the other conditions in Definition \ref{defstandardsolution}
extends to a standard solution; hence at least one standard solution exists by
short-time existence \cite[Theorem 1.1]{Shi2}.
\end{lemma}
\begin{proof}
  \uses{thm:kl-canonical-neighborhood-theorem-compactness-canonical-neighborhoods,thm:kl-no-local-collapsing-theorem-ii-existence-kappa-solutions}
\sketch (1) is the curvature-operator maximum principle \cite[Theorem 4.14]{Shi3};
(2) follows from Appendix \ref{applocalder}; (3) is the usual maximality
argument, using the derivative estimates in Appendix \ref{applocalder}; (4) is
Theorem \ref{nolocalcollapse}; and (5) is Theorem \ref{thmI.12.1}. The same
maximality argument proves the final assertion.
\end{proof}

\subsection{The blow-up time for a standard solution is $\leq 1$}
\dcref{Claim 2 of II.2} \label{claim2}

\begin{lemma}[compactness for standard solutions]
  \label{lem:kl-blow-up-time-standard-solution-compactness-standard-solutions}
  \source{kleiner-lott}
  \dcref{Claim 2 of II.2}
  \group{Kleiner--Lott Surgery Foundations}
 \label{lemclaim2}
Let $[0,T_S)$ be maximal such that the curvature of every standard solution is
uniformly bounded on each compact subinterval $[0,a]\subset[0,T_S)$. Then the
family of standard solutions converges uniformly at spatial infinity on
$[0,T_S)$ to the Ricci flow of the round infinite cylinder $S^2\times\R$ of
scalar curvature $1$. In particular, $T_S\leq1$.
\end{lemma}
\begin{proof}
  \uses{lem:kl-standard-solutions-existence-kappa-solutions,thm:kl-maximum-principles-existence-curvature-operators}
\sketch Pointed compactness from Lemma \ref{lemstandard0}(2), Appendix
\ref{applocalder}, and Appendix \ref{subsequence} reduces the assertion to a
pointed limit based at points escaping to infinity. Its time-zero slice is the
round cylinder. Each later slice has two ends and nonnegative sectional
curvature, so Toponogov splitting and the strong maximum principle split off an
$\R$ factor. Uniqueness on the round $S^2$ factor identifies the limit with the
shrinking round cylinder.
\end{proof}

\subsection{The blow-up time of a standard solution is $1$}
\dcref{Claim 4 of II.2}
\label{secclaim4}

\begin{lemma}[standard solutions lemma]
  \label{lem:kl-blow-up-time-standard-solution-standard-solutions-lemma}
  \source{kleiner-lott}
  \uses{lem:kl-blow-up-time-standard-solution-compactness-standard-solutions}
  \dcref{Claim 4 of II.2}
  \group{Kleiner--Lott Surgery Foundations}
 \label{lemclaim4}
Let $T_S$ be as in Lemma \ref{lemclaim2}. Then $T_S=1$, and every standard
solution survives until time $1$.
\end{lemma}
\begin{proof}
  \uses{lem:kl-blow-up-time-standard-solution-compactness-standard-solutions,lem:kl-canonical-neighborhood-theorem-canonical-neighborhoods-lemma-8,lem:kl-standard-solutions-existence-kappa-solutions,lem:kl-three-dimensional-solutions-existence-kappa-solutions-4,prop:kl-asymptotic-scalar-curvature-asymptotic-volume-ratio-existence-kappa-solutions,thm:kl-canonical-neighborhood-theorem-compactness-canonical-neighborhoods,thm:kl-statement-pseudolocality-theorem-pseudolocality-theorem}
\sketch Short-time existence gives $T_S>\alpha>0$. If $T_S<1$, choose
$(x_i,t_i)$ with $t_i\to T_S$ and $R(x_i,t_i)\to\infty$. The points cannot
escape to spatial infinity, since Lemma \ref{lemclaim2} and the high-curvature
derivative estimate \eqref{derestimates} give a contradiction near $T_S$.
Thus they stay in a compact region. The canonical-neighborhood theorem and
Proposition \ref{propI.11.4} give rescaled ancient limits with vanishing
asymptotic volume ratio, so Bishop--Gromov implies
$D^{-3}\vol B(x_i,t_i,D)\to0$ for every fixed $D$. This contradicts smooth
extension outside a compact set. Hence $T_S=1$.
\end{proof}

\begin{lemma}[standard solutions lemma]
  \label{lem:kl-blow-up-time-standard-solution-standard-solutions-lemma-3}
  \source{kleiner-lott}
  \dcref{Lemma 62.2}
  \group{Kleiner--Lott Surgery Foundations}
 \label{scalarblowup}
The infimal scalar curvature on the time-$t$ slice tends uniformly to infinity
over all standard solutions as $t\to1^-$.
\end{lemma}
\begin{proof}
  \uses{lem:kl-blow-up-time-standard-solution-compactness-standard-solutions,lem:kl-canonical-neighborhood-theorem-canonical-neighborhoods-lemma-8,thm:kl-canonical-neighborhood-theorem-compactness-canonical-neighborhoods}
\sketch Otherwise choose pointed standard solutions with $t_i\to1$ and bounded
$R(x_i,t_i)$. The points cannot escape to infinity by Lemma \ref{lemclaim2} and
$|(R^{-1})_t|\leq\eta$. If they remain in a compact region, bounded-distance
compactness and the canonical-neighborhood estimates produce a standard limit
defined with bounded curvature through time $1$, contradicting Lemma
\ref{lemclaim2}.
\end{proof}

\subsection{Canonical neighborhood property for standard solutions}
\dcref{Claim 5 of II.2}
\label{secclaim5}

Let $p$ be the center of the hemispherical region in the time-zero slice.

\begin{lemma}[existence for canonical neighborhoods]
  \label{lem:kl-canonical-neighborhood-property-standard-solutions-existence-canonical-neighborhoods}
  \source{kleiner-lott}
  \uses{lem:kl-three-dimensional-solutions-existence-kappa-solutions-4, lem:kl-three-dimensional-solutions-existence-kappa-solutions-5}
  \dcref{Claim 5 of II.2}
  \group{Kleiner--Lott Surgery Foundations}
 \label{claim5}
For sufficiently small $\epsilon>0$ there are constants
$\eta(\epsilon),C_1(\epsilon),C_2(\epsilon)$, independent of the standard
solution, such that the conclusions of Lemmas \ref{derivativeestimates} and
\ref{nbhd} hold, except that an $\epsilon$-neck need not be strong. Every
point $(x,t)$ falls into one of these cases:

1. If $t\in(\frac34,1)$, it has an $\epsilon$-cap or a strong
$\epsilon$-neck as in Lemma \ref{nbhd}.

2. If $x\in B(p,0,\epsilon^{-1})$ and $t\in[0,\frac34]$, it has an
$\epsilon$-cap as in Lemma \ref{nbhd}.

3. If $x\notin B(p,0,\epsilon^{-1})$ and $t\in[0,\frac34]$, there is an
$\epsilon$-neck $B(x,t,\epsilon^{-1}r)$ such that
$P(x,t,\epsilon^{-1}r,-t)$, after scaling by $r^{-2}$, is
$\epsilon$-close to the corresponding portion of the evolving round infinite
cylinder.

Moreover, for a solution-independent constant,
$R_{\min}(t)\geq\const(1-t)^{-1}$.
\end{lemma}
\begin{proof}
  \uses{lem:kl-blow-up-time-standard-solution-compactness-standard-solutions,lem:kl-blow-up-time-standard-solution-standard-solutions-lemma,lem:kl-blow-up-time-standard-solution-standard-solutions-lemma-3,lem:kl-standard-solutions-existence-kappa-solutions,lem:kl-three-dimensional-solutions-existence-kappa-solutions-4,lem:kl-three-dimensional-solutions-existence-kappa-solutions-5,thm:kl-canonical-neighborhood-theorem-compactness-canonical-neighborhoods}
\sketch Near $t=1$, apply Lemma \ref{scalarblowup} and Theorem
\ref{thmI.12.1}; away from the cap center, use the cylindrical convergence of
Lemma \ref{lemclaim2}; and in the bounded central region use the cap alternative
of Lemma \ref{nbhd}. The derivative estimates follow near $1$ from Theorem
\ref{thmI.12.1}, at initial times from Appendix \ref{applocalder}, and in the
intermediate interval from Lemma \ref{lemclaim4}, Hamilton--Ivey pinching, and
Shi estimates. The backward cylinder issue is handled by
$R(x,t)\geq(1-\frac23t)^{-1}$, which gives $r^2t>1$ for $t\geq\frac34$.
Finally $|(R^{-1})_t|\leq\const$ and $R^{-1}(x,t)\to0$ imply
$R(x,t)\geq\const(1-t)^{-1}$.
\end{proof}

\subsection{Compactness of the space of standard solutions}
\label{standardlycompact}

\begin{lemma}[compactness for standard solutions]
  \label{lem:kl-compactness-space-standard-solutions-compactness-standard-solutions}
  \source{kleiner-lott}
  \dcref{Lemma 64.1}
  \group{Kleiner--Lott Surgery Foundations}
 \label{standardcompactness}
The family ${\mathcal ST}$ of pointed standard solutions $\{(\M,(p,0))\}$ is
compact in the pointed smooth topology.
\end{lemma}
\begin{proof}
  \uses{lem:kl-blow-up-time-standard-solution-compactness-standard-solutions,lem:kl-blow-up-time-standard-solution-standard-solutions-lemma}
\sketch Apply Appendix \ref{subsequence} and Lemma \ref{lemclaim4}, which gives
$T_S=1$ in Lemma \ref{lemclaim2}.
\end{proof}

\subsection{Rotational symmetry of standard solutions}
\dcref{Claim 1 of II.2} \label{claimII.2.1}

\begin{lemma}[standard solutions lemma]
  \label{lem:kl-rotational-symmetry-standard-solutions-standard-solutions-lemma}
  \source{kleiner-lott}
  \dcref{Lemma 1}
  \group{Kleiner--Lott Surgery Foundations}
 (cf. Claim 1 of II.2)
Any Ricci flow solution in ${\mathcal ST}$ is rotationally symmetric for every
$t\in[0,1)$.
\end{lemma}
\begin{proof}
  \uses{lem:kl-blow-up-time-standard-solution-compactness-standard-solutions}
\sketch Evolve a rotational Killing field $u=\sum_m u^m\partial_m$ by
\begin{equation} \label{Killingevolution}
u_t^m=u^{m\,k}_{\ \ ;k}+R^m_{\ i}u^i.
\end{equation}
For $L_{ij}=u_{j;i}+u_{i;j}$ and $M_{ij}=e^{\lambda t}L_{ij}$, curvature
commutation and the Bianchi identity yield
\begin{align} \label{bigmess1}
\partial_tL_{ij}&=L_{ij;k}^{\ \ \ k}+2R^{\ k\ l}_{i\ j}L_{kl}-R_i^{\ k}L_{kj}-R^k_{\ j}L_{ik},
\end{align}
\begin{equation} \label{bigmess2}
R_{mlki;}^{\ k}=R_{il;m}-R_{im;l},
\end{equation}
and
\begin{align} \label{bigmess3}
\partial_t\Gamma^m_{\ ki}&=2R^{ml}\Gamma_{lki}-g^{ml}(R_{lk,i}+R_{li,k}-R_{ik,l}).
\end{align}
Consequently,
\begin{equation} \label{M2}
\partial_t(M_{ij}M^{ij})=2\lambda M_{ij}M^{ij}+(M_{ij}M^{ij})_{;k}^{\ \ k}
-2M_{ij;k}M^{ij;k}+Q(\Rm,M),
\end{equation}
where $Q$ is linear in $\Rm$ and quadratic in $M$. Choose $\lambda$ sufficiently
negative and apply the maximum principle, using $M(0)=0$ and decay at spatial
infinity, to obtain $M=0$. A Killing field satisfies
$u^{m\,k}_{\ \ ;k}+R^m_{\ i}u^i=0$, so the evolved fields are time-independent
and generate the rotational symmetry.
\end{proof}

\subsection{Uniqueness of the standard solution}
\dcref{Claim 3 of II.2} \label{claim3}

\begin{equation} \label{Weqn}
W^i(t)=h^{jk}\left(\Gamma(h)^i_{\ jk}-\Gamma(\widehat g)^i_{\ jk}\right).
\end{equation}

Two Ricci flows on $M$ with bounded curvature on compact time intervals and the
same initial metric are equal when they are standard solutions. For
$h(t)=\phi(t)^{-1*}g(t)$, where the generator of $\phi$ is $-W$, the
DeTurck equation is
\begin{equation} \label{heqn}
\frac{dh_{ij}}{dt}=-2R_{ij}(h)+\nabla(h)_iW_j+\nabla(h)_jW_i,
\end{equation}
with leading term
\begin{equation} \label{hder}
\frac{dh_{ij}}{dt}=h^{kl}\partial_k\partial_lh_{ij}+\ldots.
\end{equation}
For $w=h-\widehat g$,
\begin{equation} \label{wder}
\frac{dw}{dt}=-\nabla(\widehat g)^*\nabla(\widehat g)w+Pw+Qw,
\end{equation}
where $P$ and $Q$ are bounded first- and zeroth-order operators.
\begin{proof}
  \uses{lem:kl-blow-up-time-standard-solution-compactness-standard-solutions}
\sketch Use the harmonic heat flow to construct $\phi$ (or $F$) on a short
interval; rotational symmetry and the cylindrical asymptotics imply that it is
a diffeomorphism and converges to the identity at spatial infinity. The
DeTurck equation has principal part \eqref{hder}, and subtraction gives
\eqref{wder}. On compact exhaustions $K_i$, integrate against $w$, choose
\begin{equation}
\lambda>\sup_{v\ne0}\frac{\frac14|P^{abcdi}v_{ab}|^2+\langle v,Qv\rangle}{\langle v,v\rangle},
\end{equation}
and use the vanishing boundary terms at infinity to obtain
\begin{equation}
e^{2\lambda t}\frac{d}{dt}\left(\frac12e^{-2\lambda t}\int_{K_i}|w|^2\,d\vol_{\widehat g(t)}\right)\leq\frac1i.
\end{equation}
Since $w(0)=0$, letting $i\to\infty$ gives $w=0$, hence $W=0$ and
$g=\widehat g$. The common time interval can be extended because the equality
set is both closed and relatively open.
\end{proof}

\section{Structure at the first singularity time}
\dcref{II.3} \label{II.3}

Let $M$ be a connected closed oriented $3$-manifold and let $g(\cdot)$ be a
Ricci flow on $[0,T)$, $T<\infty$, with
$\lim_{t\to T^-}\max_{x\in M}|\Rm|(x,t)=\infty$. For sufficiently small
$\epsilon>0$, canonical-neighborhood and noncollapsing results give
$r=r(\epsilon)>0$ and $\kappa=\kappa(\epsilon)>0$ such that a point with
$Q=R(x,t)\geq r^{-2}$ has a parabolic neighborhood
$P(x,t,(\epsilon Q)^{-1/2},(\epsilon Q)^{-1})$ which is, after scaling by
$Q$, $\epsilon$-close to a $\kappa$-solution; the estimates
\eqref{derestimates} and the alternatives of Lemma \ref{nbhd} also apply.
If the sectional curvature becomes everywhere positive, $M$ is a finite round
$S^3$ quotient and shrinks to a point; hence assume this does not occur.

\begin{definition}[bounded-curvature regular set]
  \label{def:kl-structure-at-first-singularity-time-bounded-curvature-regular-set}
  \source{kleiner-lott}
  \dcref{Definition 1}
  \group{Kleiner--Lott Singularities and Surgery}

Define
\begin{equation}
\Omega=\{x\in M:\sup_{t\in[0,T)}|\Rm|(x,t)<\infty\}.
\end{equation}
\end{definition}

For $x\in M-\Omega$, the estimates \eqref{derestimates} and almost
nonnegative sectional curvature imply $R(x,t)\to\infty$ as $t\to T^-$;
thus $x\in M-\Omega$ exactly when this scalar-curvature limit is infinite.

\begin{lemma}[structure at the first singularity time lemma]
  \label{lem:kl-structure-at-first-singularity-time-structure-at-first-singularity-time-lemma}
  \source{kleiner-lott}
  \dcref{Lemma 67.3}
  \group{Kleiner--Lott Singularities and Surgery}
 \label{omegaopen}
$\Omega$ is open in $M$.
\end{lemma}
\begin{proof}
\sketch For $x\in\Omega$, the inverse-scalar estimates
\eqref{eqnderestimatesrinverse} bound $R$ on a uniform spatial neighborhood of
$x$ for all $t<T$. Almost nonnegative sectional curvature bounds $|\Rm|$ there,
and length distortion converts this to a fixed neighborhood in $M$ contained
in $\Omega$.
\end{proof}

\begin{lemma}[compactness for structure at the first singularity time]
  \label{lem:kl-structure-at-first-singularity-time-compactness-structure-at-first-singularity-time}
  \source{kleiner-lott}
  \dcref{Lemma 67.4}
  \group{Kleiner--Lott Singularities and Surgery}
 \label{omeganoncompact}
Every connected component $C$ of $\Omega$ is noncompact.
\end{lemma}
\begin{proof}
\sketch A compact component would be all of the connected manifold $M$, which
contradicts the singularity at time $T$.
\end{proof}

\begin{lemma}[rigidity for structure at the first singularity time]
  \label{lem:kl-structure-at-first-singularity-time-rigidity-structure-at-first-singularity-time}
  \source{kleiner-lott}
  \dcref{Lemma 67.5}
  \group{Kleiner--Lott Singularities and Surgery}
 \label{standard}
If $\Omega=\emptyset$, then
$M\cong S^3$, $\R P^3$, $S^1\times S^2$, or
$\R P^3\#\R P^3$.
\end{lemma}
\begin{proof}
  \uses{lem:kl-three-dimensional-solutions-existence-kappa-solutions,lem:kl-three-dimensional-solutions-existence-kappa-solutions-5}
\sketch Near $T$, every point lies in an $\epsilon$-neck or $\epsilon$-cap.
Splicing neck projections gives an $S^2$-fibration over a $1$-manifold. If
caps occur, there are exactly two and capping the product region gives
$S^3$, $\R P^3$, or $\R P^3\#\R P^3$; if no caps occur, orientability makes the
$S^2$-bundle over $S^1$ equal to $S^1\times S^2$.
\end{proof}

Assume henceforth that $\Omega\ne\emptyset$. The local derivative estimates
give a smooth limit metric $\overline g=\lim_{t\to T^-}g(t)|_\Omega$ with scalar
curvature $\overline R$.

\begin{lemma}[volume lemma]
  \label{lem:kl-structure-at-first-singularity-time-volume-lemma}
  \source{kleiner-lott}
  \dcref{Lemma 67.6}
  \group{Kleiner--Lott Singularities and Surgery}
 \label{volumeomega}
$(\Omega,\overline g)$ has finite volume.
\end{lemma}
\begin{proof}
\sketch Integrating
$\frac{d}{dt}\vol(M,g(t))=-\int_MR\,d\vol_M$ and using
\eqref{Rlowerbound} gives $\vol(M,g(t))\leq\const+\const t^{3/2}$ for
$t<T$, hence the limit region has finite volume.
\end{proof}

\begin{lemma}[existence for structure at the first singularity time]
  \label{lem:kl-structure-at-first-singularity-time-existence-structure-at-first-singularity-time}
  \source{kleiner-lott}
  \dcref{Lemma 67.7}
  \group{Kleiner--Lott Singularities and Surgery}

\label{lemrinverseextends}
There is an open neighborhood $V$ of $(M-\Omega)\times\{T\}$ in
$M\times[0,T]$ on which $R^{-1}$ extends continuously and vanishes on
$(M-\Omega)\times\{T\}$.
\end{lemma}
\begin{proof}
  \uses{lem:kl-structure-at-first-singularity-time-structure-at-first-singularity-time-lemma}
\sketch The characterization of $M-\Omega$ by the limit of $R^{-1}$ and the
spacetime derivative estimates \eqref{eqnderestimatesrinverse} give the
continuous extension along suitable spacetime paths.
\end{proof}

\begin{definition}[curvature sublevel of the regular set]
  \label{def:kl-structure-at-first-singularity-time-curvature-sublevel-regular-set}
  \source{kleiner-lott}
  \dcref{Definition 7}
  \group{Kleiner--Lott Singularities and Surgery}

For $\rho<\frac r2$, define
$\displaystyle\Omega_\rho=\{x\in\Omega:\overline R(x)\leq\rho^{-2}\}$.
\end{definition}

\begin{lemma}[characterization for structure at the first singularity time]
  \label{lem:kl-structure-at-first-singularity-time-characterization-structure-at-first-singularity-time}
  \source{kleiner-lott}
  \dcref{Lemma 67.9}
  \group{Kleiner--Lott Singularities and Surgery}

\label{lemoverlinerisproper}
The function $\overline R:\Omega\to\R$ is proper. Equivalently, if
$x_i\in\Omega$ leaves every compact subset of $\Omega$, then
$\overline R(x_i)\to\infty$. In particular, $\Omega_\rho$ is compact in $M$
for every $\rho<r$.
\end{lemma}
\begin{proof}
\sketch If $\overline R(x_i)$ were bounded, a subsequence would converge in
$M$. The continuity and vanishing statement for $R^{-1}$ from Lemma
\ref{lemrinverseextends} force the limit to lie in $\Omega$, a contradiction.
\end{proof}

For a component $C$ of $\Omega$ disjoint from $\Omega_\rho$, use the following
terminology. In $\R^3$ or $\R P^3-\overline{B^3}$, an embedded $2$-sphere
separates a compact side from a noncompact side. A good cylinder is an open
region which is, modulo rescaling, $\epsilon$-close to a round cylindrical
segment, has diameter approximately $100$ times its cross-section, and whose
points all lie in ambient $\epsilon$-necks. Every $\epsilon$-cap contains such
a good cylinder in its end neck.

\begin{lemma}[neck geometry lemma]
  \label{lem:kl-structure-at-first-singularity-time-neck-geometry-lemma}
  \source{kleiner-lott}
  \uses{lem:kl-three-dimensional-solutions-existence-kappa-solutions-5}
  \dcref{Lemma 67.10}
  \group{Kleiner--Lott Singularities and Surgery}
 \label{doublehorn}
If every $x\in C$ has a neighborhood $B_x$ of the strong $\epsilon$-neck type
in Lemma \ref{nbhd}, then $C$ is a double $\epsilon$-horn.
\end{lemma}
\begin{proof}
  \uses{lem:kl-structure-at-first-singularity-time-existence-structure-at-first-singularity-time,lem:kl-three-dimensional-solutions-existence-kappa-solutions}
\sketch Glue the overlapping necks, using good cylinders with successive
overlaps of cross-sectional diameter approximately $10$, to obtain an
$S^2$-fibration over a $1$-manifold. Lemma \ref{lemrinverseextends} forces both
ends to have unbounded scalar curvature.
\end{proof}

\begin{lemma}[existence for cap geometry]
  \label{lem:kl-structure-at-first-singularity-time-existence-cap-geometry}
  \source{kleiner-lott}
  \uses{lem:kl-three-dimensional-solutions-existence-kappa-solutions-5}
  \dcref{Lemma 10}
  \group{Kleiner--Lott Singularities and Surgery}

If some $x\in C$ has an $\epsilon$-cap neighborhood $B_x$ as in Lemma
\ref{nbhd}, then $C$ is a capped $\epsilon$-horn.
\end{lemma}
\begin{proof}
  \uses{lem:kl-structure-at-first-singularity-time-compactness-structure-at-first-singularity-time,lem:kl-structure-at-first-singularity-time-volume-lemma,lem:kl-three-dimensional-solutions-existence-kappa-solutions-5}
\sketch Starting from the cap, glue successive good cylinders. If the process
continues indefinitely, the resulting open set is also closed by finite volume,
so it equals $C$. Otherwise each obstruction is another cap, producing nested
caps and disjoint good cylinders. The estimate
$\vol(B_{p_k})\geq\const k\,\overline R(p_1)^{-3/2}$ contradicts finite volume
for large $k$; the finite process therefore ends in a capped horn. In the
separation step, if $V,V'$ are the compact sides of the same cross-sectional
sphere, then
\begin{equation}
\label{eqndecompv}
V=(V\cap V')\sqcup(V-V'),\qquad V'=(V\cap V')\sqcup(V'-V),
\end{equation}
and connectedness gives $V=V'$.
\end{proof}

There may be infinitely many components of $\Omega$ disjoint from $\Omega_\rho$.
For a component meeting $\Omega_\rho$, the neck-gluing process terminates at
$\Omega_{2\rho}$ in at least one direction. For sufficiently small $\epsilon$,
each component of $C-(C\cap\Omega_\rho)$ is an $\epsilon$-tube with both
boundaries in $\Omega_\rho$, an $\epsilon$-cap with boundary in $\Omega_\rho$,
or an $\epsilon$-horn with boundary in $\Omega_\rho$. Quantitative volume
lower bounds and Lemma \ref{volumeomega} imply that only finitely many
components of $\Omega$ meet $\Omega_\rho$ and that each has finitely many
disjoint $\epsilon$-horn ends.

For each such component, cut far-out cross-sectional spheres in its horns,
discard the tips and components disjoint from $\Omega_\rho$, and cap the
remaining boundary spheres by $3$-balls to obtain a possibly disconnected
postsurgery manifold $M'$.

\begin{lemma}[structure at the first singularity time lemma]
  \label{lem:kl-structure-at-first-singularity-time-structure-at-first-singularity-time-lemma-11}
  \source{kleiner-lott}
  \dcref{Lemma 67.13}
  \group{Kleiner--Lott Singularities and Surgery}
 \label{reconstruct}
$M$ is obtained from components of $M'$ by connected sums, with possibly a
finite number of additional connected sums with $S^1\times S^2$ and $\R P^3$.
\end{lemma}
\begin{proof}
  \uses{lem:kl-three-dimensional-solutions-existence-kappa-solutions-5}
\sketch Components of $M-X$ are $\epsilon$-tubes or $\epsilon$-caps. A tube
whose boundary spheres lie in distinct components of $X$ gives a connected sum
of those components; a tube attached to one component gives an
$S^1\times S^2$ summand. A $3$-ball cap has no effect, while an
$(\R P^3-B^3)$ cap gives an $\R P^3$ summand.
\end{proof}

\begin{remark}[estimates for diameter]
  \label{rem:kl-structure-at-first-singularity-time-estimates-diameter}
  \source{kleiner-lott}
  \dcref{Remark 12}
  \group{Kleiner--Lott Singularities and Surgery}

No boundedness of $\diam(M,g(t))$ as $t\to T$ is assumed.
\end{remark}

\section{Ricci flow with surgery: the general setting}
\label{Ricci flow with surgery}

We use the notation $\M$, $\M_t$, and $\M_{\reg}$ for a Ricci flow with
surgery, its time-$t$ slice, and its regular points. At a singular time,
$\M_t^-$ is the limiting presurgery slice and $\M_t^+$ is the outgoing
postsurgery slice; at a nonsingular time they agree.

\begin{definition}[ricci flow with surgery definition]
  \label{def:kl-ricci-flow-surgery-general-setting-ricci-flow-surgery-definition}
  \source{kleiner-lott}
  \dcref{Definition 68.1}
  \group{Kleiner--Lott Singularities and Surgery}
  \label{flowwithsurgery}
A Ricci flow with surgery consists of:

\begin{itemize}
\item Ricci flows $\{(M_k\times[t_k^-,t_k^+),g_k) \}_{1\leq k\leq N}$,
where $N\leq\infty$, the $M_k$ are compact (possibly empty),
$t_k^+=t_{k+1}^-$ for $k<N$, and $g_k$ becomes singular at $t_k^+$ for
$k<N$; $t_N^+$ may be infinite.
\item Limits $(\Om_k,\bar g_k)$ at the singular final times, with
$\Omega_k$ open in $M_k$.
\item Isometric embeddings
$\psi_k:X_k^+\ra X_{k+1}^-$ for compact $3$-submanifolds with boundary
$X_k^+\subset\Om_k$ and $X_{k+1}^-\subset M_{k+1}$.
\end{itemize}
The subsets $X_k^\pm$ survive the transition and the $\psi_k$ identify
them. A time $t$ is singular if $t=t_k^+=t_{k+1}^-$ for some $k<N$, or
if $t=t_N^+$ and the final metric becomes singular.
\end{definition}

The associated spacetime is obtained by gluing
\begin{equation}
\left(M_k\times [t_k^-,t_k^+)\right)\cup
\left(\Om_k\times\{t_k^+\}\right)
\subset M_k\times[t_k^-,t_k^+]
\end{equation}
along the $\psi_k$. Its quotient map is $\pi$ and its time slices are
denoted by $\M_t$. At $t=t_k^+=t_{k+1}^-$,
$\M_t^- = \pi(\Omega_k\times\{t_k^+\})$ and
$\M_t^+ = \pi(M_{k+1}\times\{t_{k+1}^-\})$; the glued slice consists of
the surviving regions and their identification. Set
$\M_{(t,t')}=\bigcup_{\bar t\in(t,t')}\M_{\bar t}$ and
$\M_{[t,t']}=\M_t^+\cup\M_{(t,t')}\cup\M_{t'}^-$. The spacetime is a
topological manifold with boundary, with regular points the smooth
$4$-manifold points together with the initial and final slices; splitting
points are images of $\D X_k^+\times\{t_k^+\}$. The horizontal metric is
smooth on $\M_{\reg}$.

Distances on $\M_t$ and $\M_t^\pm$ are induced by piecewise smooth path
length; $B(x,t,r)$ and $B^\pm(x,t,r)$ denote the corresponding balls.
Proper balls are those for which the distance function to the center is
proper. An admissible curve $\ga:[c,d]\ra\M$ is time-preserving and lifts
smoothly through each constituent flow; $\dot\ga$ denotes its horizontal
velocity. Static curves have constant spatial component. A barely admissible
curve meeting a surgery region passes through
\begin{equation}
\pi(\D X_k^+\times\{t_k^+\})
=\pi(\D X_{k+1}^-\times\{t_{k+1}^-\}).
\end{equation}
The parabolic region $P(x,t,r,\De t)$ is the union of static curves from
$B^+(x,t,r)$ for $\De t>0$, or ending in $B^-(x,t,r)$ for $\De t<0$.
A set is unscathed on $[c,d]$ if each point lies on a static curve defined
through that interval. A strong $\eps$-neck is an unscathed product
$U\times[c,d]$ whose terminal slice is a strong $\eps$-neck.

\subsection{\emph{A priori} assumptions}
\dcref{II.4.1}

\begin{definition}[canonical neighborhood assumption]
  \label{def:kl-priori-assumptions-canonical-neighborhood-assumption}
  \source{kleiner-lott}
  \uses{lem:kl-canonical-neighborhood-property-standard-solutions-existence-canonical-neighborhoods, lem:kl-three-dimensional-solutions-existence-kappa-solutions-4, lem:kl-three-dimensional-solutions-existence-kappa-solutions-5}
  \dcref{Definition II.2}
  \group{Kleiner--Lott Singularities and Surgery}
\label{canonicalnbhddef}
Choose $\eps>0$ so that Lemmas \ref{nbhd} and \ref{claim5} hold. Let
$C_1$ be the maximum of $30\epsilon^{-1}$ and their $C_1(\eps)$ values, and
let $C_2$ be the maximum of their $C_2(\epsilon)$ values. For a positive
nonincreasing $r:[a,b]\ra(0,\infty)$, a Ricci flow with surgery on $[a,b]$
satisfies the $r$-canonical neighborhood assumption if every
$(x,t)\in\M_t^\pm$ with $R(x,t)\geq r(t)^{-2}$ has an open canonical
neighborhood $U$ with
\[
\ol{B^\pm(x,t,\hat r)}\subset U\subset B^\pm(x,t,2\hat r),\qquad
R(x,t)^{-1/2}<\hat r<C_1R(x,t)^{-1/2},
\]
and one of:
\begin{enumerate}
\item $U\times[t-\De t,t]$ is a strong $\eps$-neck;
\item $U$ is an $\eps$-cap $\epsilon$-close after rescaling to a piece of a
$\kappa_0$-solution or a time slice of a standard solution;
\item $U$ is a closed manifold diffeomorphic to $S^3$ or $RP^3$;
\item $U$ is $\epsilon$-close to a closed manifold of constant positive
sectional curvature.
\end{enumerate}
In all cases $C_2^{-1}R(x,t)\leq R|_U\leq C_2R(x,t)$. In cases (1)--(3),
$\vol(U)>C_2^{-1}R(x,t)^{-3/2}$, and in case (3) the infimal sectional
curvature is greater than $C_2^{-1}R(x,t)$. Finally,
\begin{equation}
\label{II.(1.3)estimates}
|\nabla R(x,t)|<\eta R(x,t)^{3/2},\qquad
\left|\frac{\D R}{\D t}(x,t)\right|<\eta R(x,t)^2,
\end{equation}
where the time derivative is one-sided at surgery points and $\eta$ is the
constant from (\ref{derestimates}).
\end{definition}

\begin{remark}[continuation for canonical neighborhoods]
  \label{rem:kl-priori-assumptions-continuation-canonical-neighborhoods}
  \source{kleiner-lott}
  \uses{lem:kl-performing-surgery-continuing-flows-existence-canonical-neighborhoods, rem:kl-performing-surgery-continuing-flows-structural-remark-canonical-neighborhoods}
  \dcref{Remark 2}
  \group{Kleiner--Lott Singularities and Surgery}
The closed inner ball in the canonical-neighborhood condition facilitates
openness. The lower bound $C_1\geq30\epsilon^{-1}$ is used in Lemma
\ref{extendit}; see Remark \ref{30}.
\end{remark}

\begin{remark}[continuation for kappa solutions]
  \label{rem:kl-priori-assumptions-continuation-kappa-solutions}
  \source{kleiner-lott}
  \dcref{Remark 3}
  \group{Kleiner--Lott Singularities and Surgery}
For convenience, case (b) includes $\epsilon$-closeness to a
$\kappa_0$-solution or standard solution. Added points lie in $\eps$-caps,
while points farther out on the capped neck lie in strong $\eps$-necks
extending backward before surgery.
\end{remark}

\begin{definition}[curvature pinching assumption]
  \label{def:kl-priori-assumptions-curvature-pinching-assumption}
  \source{kleiner-lott}
  \lean{KleinerLott.IsPinchingFunction, KleinerLott.IsCurvaturePinchedBy}
  \dcref{Definition 69.5}
  \group{Kleiner--Lott Singularities and Surgery}
\label{phipinchingdef}
($\Phi$-pinching) Let $\Phi\in C^\infty(\R)$ be positive and nondecreasing,
with $\Phi(s)/s$ decreasing for positive $s$ and tending to zero as
$s\ra\infty$. The $\Phi$-pinching assumption is
\[
\Rm(x,t)\geq-\Phi(R(x,t))\quad\text{for all }(x,t)\in\M.
\]
\end{definition}

\begin{definition}[surgery a priori assumptions]
  \label{def:kl-priori-assumptions-surgery-priori-assumptions}
  \source{kleiner-lott}
  \uses{def:kl-priori-assumptions-canonical-neighborhood-assumption, def:kl-priori-assumptions-curvature-pinching-assumption}
  \dcref{Definition 69.6}
  \group{Kleiner--Lott Singularities and Surgery}
 \label{apriori}
A Ricci flow with surgery satisfies the \emph{a priori} assumptions if it
satisfies the $\Phi$-pinching and $r$-canonical neighborhood assumptions.
They depend on $\epsilon$, $r(t)$ from Definition \ref{canonicalnbhddef},
and $\Phi$ from Definition \ref{phipinchingdef}.
\end{definition}

\subsection{Curvature bounds from the \emph{a priori} assumptions}
\dcref{II.4.2}

\begin{lemma}[estimates for curvature]
  \label{lem:kl-curvature-bounds-from-priori-assumptions-estimates-curvature}
  \source{kleiner-lott}
  \uses{def:kl-priori-assumptions-canonical-neighborhood-assumption}
  \dcref{Claim 1 of II.4.2}
  \group{Kleiner--Lott Singularities and Surgery}
\label{claim1II.4.2}
Given $(x_0,t_0)\in\M$, set $Q=|R(x_0,t_0)|+r(t_0)^{-2}$. Then
\[
R(x,t)\leq8Q\quad\text{for }(x,t)\in
P\left(x_0,t_0,\tfrac12\eta^{-1}Q^{-1/2},-\tfrac18\eta^{-1}Q^{-1}\right),
\]
where $\eta$ is from (\ref{II.(1.3)estimates}).
\end{lemma}
\begin{proof}
  \uses{def:kl-priori-assumptions-canonical-neighborhood-assumption,lem:kl-canonical-neighborhood-theorem-canonical-neighborhoods-lemma-8}
\sketch Integrate the spatial and temporal derivative estimates along a
spatial path followed by a static backward path; monotonicity of $r$ ensures
the estimates apply at all relevant points.
\end{proof}

\begin{lemma}[canonical neighborhoods lemma]
  \label{lem:kl-curvature-bounds-from-priori-assumptions-canonical-neighborhoods-lemma}
  \source{kleiner-lott}
  \dcref{Claim 2 of II.4.2}
  \group{Kleiner--Lott Singularities and Surgery}
\label{claim2II.4.2}
If $\epsilon$ is sufficiently small and $\M$ satisfies the $\Phi$-pinching
and $r$-canonical neighborhood assumptions, then for every $A<\infty$ and
$\widehat r>0$ there are $\xi=\xi(A)>0$ and $K=K(A,\widehat r)<\infty$ such
that, whenever $Q=R(x_0,t_0)>0$, $\Phi(Q)/Q<\xi$, and
$\dist_{t_0}(x_0,x)\leq A Q^{-1/2}$, one has
$R(x,t_0)\leq KQ$, with $K=K(A,r(t_0))$.
\end{lemma}
\begin{proof}
  \uses{def:kl-priori-assumptions-canonical-neighborhood-assumption,thm:kl-canonical-neighborhood-theorem-compactness-canonical-neighborhoods}
\sketch A failed bound produces a nonflat metric-cone blowup. Canonical
neighborhoods force the blowup points into strong $\epsilon$-necks, where the
strong maximum principle contradicts the cone limit.
\end{proof}

\subsection{$\delta$-necks in $\epsilon$-horns}
\dcref{II.4.3}

\begin{lemma}[existence for Ricci flow with surgery]
  \label{lem:kl-necks-horns-existence-ricci-flow-surgery}
  \source{kleiner-lott}
  \uses{def:kl-priori-assumptions-surgery-priori-assumptions}
  \dcref{Lemma II.4.3}
  \group{Kleiner--Lott Singularities and Surgery}
(cf. Lemma II.4.3) \label{lemmaII.4.3}
\label{hexists}
Given $\Phi$, $\widehat T>0$, a positive nonincreasing
$r:[0,\widehat T]\ra\R$, and $\delta\in(0,1/2)$, there is a positive
nonincreasing $h:[0,\widehat T]\ra\R$ with $0<h(T)<\delta^2r(T)$ such that:
if $\M$ is a Ricci flow with surgery on $[0,T)$, $T<\widehat T$, satisfying
the \emph{a priori} assumptions and becoming singular at $T$, let
$(\Omega,\overline g)$ be its time-$T$ limit and put
\[
\Om_\rho=\{(x,T)\in\Omega\mid\overline R(x,T)\leq\rho^{-2}\},\qquad
\rho=\de r(T).
\]
If $(x,T)$ lies in an $\eps$-horn with boundary in $\Om_\rho$ and
$\overline R(x,T)\geq h(T)^{-2}$, then
\[
P\left(x,T,\de^{-1}\overline R(x,T)^{-1/2},-\overline R(x,T)^{-1}\right)
\]
is contained in a strong $\de$-neck.
\end{lemma}
\begin{proof}
  \uses{def:kl-priori-assumptions-canonical-neighborhood-assumption,lem:kl-blow-up-time-standard-solution-compactness-standard-solutions,lem:kl-curvature-bounds-from-priori-assumptions-canonical-neighborhoods-lemma,lem:kl-structure-at-first-singularity-time-characterization-structure-at-first-singularity-time}
\sketch If the conclusion failed, rescale at points deep in increasingly long
horns. The curvature and injectivity estimates give a complete nonnegatively
curved limit containing a bi-infinite minimizing geodesic, hence an
$\R$-splitting. The resulting ancient product is the shrinking cylinder,
which forces the prescribed parabolic regions to be strong $\delta$-necks.
\end{proof}

\begin{remark}[monotonicity for neck geometry]
  \label{rem:kl-necks-horns-monotonicity-neck-geometry}
  \source{kleiner-lott}
  \uses{lem:kl-necks-horns-existence-ricci-flow-surgery}
  \dcref{Remark 2}
  \group{Kleiner--Lott Singularities and Surgery}
The function $h$ may be chosen to depend only on $\min r=r(T)$ and to be
nondecreasing in both $r(T)$ and $\delta$.
\end{remark}

\subsection{Surgery and the pinching condition} \label{surgery}

\begin{lemma}[Ricci flow with surgery lemma]
  \label{lem:kl-surgery-pinching-condition-ricci-flow-surgery-lemma}
  \source{kleiner-lott}
  \dcref{Lemma 72.1}
  \group{Kleiner--Lott Singularities and Surgery}
\label{lemroundoff}
Let $g_{cyl}$ be the round cylindrical metric of scalar curvature $1$ on
$\R\times S^2$. Given $A>0$, let $f:(-A,0]\ra\R$ be smooth with
$f^{(k)}(0)=0$ for all $k$, $f<0$, $f'>0$, $f''<0$ on $(-A,0)$,
$\|f\|_{C^2}<\eps$, and
\begin{equation}
\label{eqnderivcomparision}
\max\{|f(z)|,|f'(z)|\}\leq\eps|f''(z)|.
\end{equation}
If $\|h_0-g_{cyl}\|_{C^2}<\eps$ on $(-A,0]\times S^2$ and
$h_1=e^{2f(z)}h_0$, then
\[
R_{h_1}(p)>R_{h_0}(p)-f''(z(p)),\qquad
\lambda_1^{h_1}(p)>\lambda_1^{h_0}(p)-f''(z(p)).
\]
\end{lemma}
\begin{proof}
\sketch The conformal scalar-curvature and curvature-operator formulas below,
combined with the cylinder closeness and derivative comparison, give the
claimed strict improvements.
\begin{equation} \label{Rvar}
\lambda_1(p)=\inf_{\omega\ne0}
\frac{\omega_{ij}R^{ij}_{\ \ kl}\omega^{kl}}{\omega_{ij}\omega^{ij}}.
\end{equation}
\begin{equation} \label{confscalar}
R_{h_1}=e^{-2f}(R_{h_0}-4\triangle f-2|\nabla f|^2).
\end{equation}
\begin{equation} \label{confcurv}
R^{ij}_{\ \ kl}(h_1)=e^{-2f}\bigl(R^{ij}_{\ \ kl}(h_0)
-\widetilde f^i_{\ k}\delta^j_{\ l}+\widetilde f^i_{\ l}\delta^j_{\ k}
+\widetilde f^j_{\ k}\delta^i_{\ l}-\widetilde f^j_{\ l}\delta^i_{\ k}
-|\nabla f|^2(\delta^i_{\ k}\delta^j_{\ l}-\delta^i_{\ l}\delta^j_{\ k})\bigr),
\end{equation}
where $\widetilde f=\Hess(f)-df\otimes df$. Applying these formulas,
the cylinder closeness, and (\ref{eqnderivcomparision}) gives the two
strict curvature improvements.
\begin{equation} \label{min1}
\lambda_1^{h_1}(p)=e^{-2f(z)}\left(\inf_{\omega\ne0}
\frac{\omega_{ij}R^{ij}_{\ \ kl}(h_0)\omega^{kl}
-4\omega_{ij}\widetilde f^i_{\ k}\omega^{kj}}{\omega_{ij}\omega^{ij}}
-2|\nabla f|^2(z)\right).
\end{equation}
\end{proof}

\begin{lemma}[positivity for Ricci flow with surgery]
  \label{lem:kl-surgery-pinching-condition-positivity-ricci-flow-surgery}
  \source{kleiner-lott}
  \uses{lem:kl-surgery-pinching-condition-ricci-flow-surgery-lemma}
  \dcref{Lemma 72.20}
  \group{Kleiner--Lott Singularities and Surgery}
 \label{capoff}
For every $A>0$ there are $B>A$ and $F\in C^\infty(-B,\infty)$ such that
\begin{enumerate}
\item $F\equiv0$ on $[0,\infty)$;
\item $F|_{(-A,0]}$ satisfies the hypotheses of Lemma \ref{lemroundoff};
\item $e^{2F(z)}g_{cyl}$ has nonnegative sectional curvature and extends
smoothly across $\{-B\}\times S^2$ by a constant positively curved ball.
\end{enumerate}
\end{lemma}
\begin{proof}
  \uses{lem:kl-surgery-pinching-condition-ricci-flow-surgery-lemma}
\sketch For $e^{2F}g_{cyl}$ the sectional curvatures are
$-e^{-2F}F''$ and $e^{-2F}(1/2-(F')^2)$. Choose $F$ satisfying the
roundoff hypotheses and interpolate its derivative to the derivative of a
round $3$-sphere metric, then glue the corresponding constant-curvature ball.
\end{proof}

\begin{lemma}[existence for standard solutions]
  \label{lem:kl-surgery-pinching-condition-existence-standard-solutions}
  \source{kleiner-lott}
  \uses{def:kl-almost-nonnegative-curvature-positivity-curvature-operators}
  \dcref{Lemma 72.24}
  \group{Kleiner--Lott Singularities and Surgery}
\label{preservingphipinching}
There exist $\de_0>0$ and $\de'=\de'(\de)$ with
$\lim_{\de\ra0}\de'(\de)=0$ such that, for $\de<\de_0$, if
$x\in\{0\}\times S^2$ and $h_0$ on $(-A,1/\de)\times S^2$ satisfies the
time-$t$ Hamilton-Ivey pinching condition, $R(x)>0$, and
$R(x)h_0$ is $\de$-close to $g_{cyl}$ in $C^{[1/\de]+1}$, then there is a
smooth metric $h$ on $\R^3=D^3\cup((-B,1/\de)\times S^2)$ such that:
\begin{itemize}
\item $h$ satisfies the time-$t$ pinching condition;
\item $h=h_0$ on $[0,1/\de)\times S^2$;
\item $R(x)h=g_0$ on $(-B,-A)\times S^2$, where $g_0$ is the standard
solution initial metric;
\item $R(x)h$ has constant curvature $k^2$ on $D^3$;
\item $R(x)h$ is $\de'$-close to $e^{2F}g_{cyl}$ in
$C^{[1/\de']+1}$ on $(-B,1/\delta)\times S^2$.
\end{itemize}
\end{lemma}

\subsection{Performing surgery and continuing flows}
\dcref{II.4.4}
\label{II.4.4}

\begin{definition}[Ricci flow with cutoff]
  \label{def:kl-performing-surgery-continuing-flows-ricci-flow-cutoff}
  \source{kleiner-lott}
  \uses{def:kl-priori-assumptions-surgery-priori-assumptions, lem:kl-necks-horns-existence-ricci-flow-surgery, lem:kl-surgery-pinching-condition-existence-standard-solutions}
  \dcref{Definition 73.1}
  \group{Kleiner--Lott Singularities and Surgery}
\label{rdeltacutoffflowdef}
Let $\M$ be a Ricci flow with surgery on $[a,b]$ satisfying the
\emph{a priori} assumptions, and let $\de:[a,b]\ra(0,\delta_0)$ be nonincreasing,
where $\delta_0$ is from Lemma \ref{preservingphipinching}. It has
$(r,\de)$-cutoff if at each singular time $t$, with
$\Omega=\M_t^-$, the outgoing slice is obtained by:
\begin{enumerate}
\item discarding components not meeting
\[
\Om_\rho=\{(x,t)\in\Omega\mid R(x,t)\leq\rho^{-2}\},\qquad
\rho=\de(t)r(t);
\]
\item in each remaining $\eps$-horn $\h_{ij}$, choosing
$R(x_{ij},t)=h(t)^{-2}$;
\item choosing the strong $\de$-neck
$U_{ij}\times[t-h(t)^2,t]$ containing
$P(x_{ij},t,\de^{-1}R(x_{ij},t)^{-1/2},-R(x_{ij},t)^{-1})$;
\item cutting along cross-sectional spheres $S_{ij}\subset U_{ij}$ and
discarding the horn tips, producing a compact manifold $X$;
\item gluing caps to $X$ by Lemma \ref{preservingphipinching}.
\end{enumerate}
\end{definition}

Take $A=10$ in the cap construction. The boundary neighborhood is
$[-A,\de^{-1})\times S^2$, the unchanged region is
$[0,\de^{-1})\times S^2$, and the added region is a $3$-disk bounded by
$S'_{ij}=\{0\}\times S^2$. In the notation of Definition
\ref{flowwithsurgery}, $\partial X_k^+=\bigcup_{ij}S'_{ij}$ and the
added part is a union of $3$-balls.

\begin{remark}[existence for canonical neighborhoods]
  \label{rem:kl-performing-surgery-continuing-flows-existence-canonical-neighborhoods}
  \source{kleiner-lott}
  \dcref{Remark 2}
  \group{Kleiner--Lott Singularities and Surgery}
The definition retains postsurgery components close to round spherical
quotients and components with nonnegative scalar curvature; omitting them as
in some alternative surgery conventions does not affect the canonical
neighborhood framework.
\end{remark}

\begin{lemma}[rigidity for Ricci flow with surgery]
  \label{lem:kl-performing-surgery-continuing-flows-rigidity-ricci-flow-surgery}
  \source{kleiner-lott}
  \dcref{Lemma 73.4}
  \group{Kleiner--Lott Singularities and Surgery}
 \label{reconstruct2}
The presurgery manifold is obtained from the postsurgery manifold by finitely
many applications of:
\begin{itemize}
\item replacing two connected components by their connected sum;
\item taking connected sum with $S^1\times S^2$ or $\R P^3$;
\item adjoining $S^1\times S^2$ or an isometric quotient of round $S^3$.
\end{itemize}
\end{lemma}

\begin{remark}[structural remark on Ricci flow with surgery]
  \label{rem:kl-performing-surgery-continuing-flows-structural-remark-ricci-flow-surgery}
  \source{kleiner-lott}
  \dcref{Remark 73.5}
  \group{Kleiner--Lott Singularities and Surgery}
\label{volumelossremark}
For sufficiently small $\de$, every surgery time satisfies
\[
\vol(\M_t^+)<\vol(\M_t^-)-h(t)^3.
\]
\end{remark}

\begin{remark}[existence for Ricci flow with surgery]
  \label{rem:kl-performing-surgery-continuing-flows-existence-ricci-flow-surgery}
  \source{kleiner-lott}
  \dcref{Remark 5}
  \group{Kleiner--Lott Singularities and Surgery}
For a nonaspherical irreducible initial manifold, surgery preserves a
component homotopy-equivalent to the presurgery manifold while other
components are spherical. A homotopy equivalence from nearby presurgery
slices to the surviving component can be chosen with distance expansion
$\xi(t')\ra1$ as $t'\ra t$; the same control applies at later singular times.
\end{remark}

\begin{lemma}[existence for canonical neighborhoods]
  \label{lem:kl-performing-surgery-continuing-flows-existence-canonical-neighborhoods}
  \source{kleiner-lott}
  \uses{def:kl-almost-nonnegative-curvature-positivity-curvature-operators}
  \dcref{Lemma 73.7}
  \group{Kleiner--Lott Singularities and Surgery}
(Prolongation of Ricci flows with cutoff)
\label{extendit}
Let $\Phi$ be the time-dependent pinching function from Definition
\ref{HamIvey}. If $r$ and $\de$ are positive nonincreasing functions on
$[a,b]$, and $\M$ is a Ricci flow with $(r,\de)$-cutoff on $[a,c]$, then for
sufficiently small $\sup\de$, either:
\begin{enumerate}
\item $\M$ prolongs to a Ricci flow with $(r,\de)$-cutoff on $[a,b]$; or
\item $\M$ extends to a Ricci flow with surgery on $[a,T]$ for some
$T\in(c,b]$, every restriction to $[a,T']$, $T'<T$, has $(r,\de)$-cutoff,
and the $r$-canonical neighborhood assumption fails at some
$(x,T)\in\M_T^-$.
\end{enumerate}
Thus failure of the canonical neighborhood assumption is the only
prolongation obstruction.
\end{lemma}
\begin{proof}
  \uses{def:kl-almost-nonnegative-curvature-positivity-curvature-operators,def:kl-performing-surgery-continuing-flows-ricci-flow-cutoff,def:kl-priori-assumptions-canonical-neighborhood-assumption,lem:kl-canonical-neighborhood-property-standard-solutions-existence-canonical-neighborhoods,lem:kl-surgery-pinching-condition-existence-standard-solutions}
\sketch Perform the cutoff surgery at the terminal slice. The cap lemma
preserves pinching and canonical neighborhoods near the added regions, while
the presurgery neighborhoods persist elsewhere. Continue the flow until it
either reaches $b$ or the canonical-neighborhood assumption fails. If
singular times accumulated, each surgery would remove at least $h^3$ of
volume, contradicting the volume bound from the maximum principle.
\end{proof}

\begin{remark}[structural remark on canonical neighborhoods]
  \label{rem:kl-performing-surgery-continuing-flows-structural-remark-canonical-neighborhoods}
  \source{kleiner-lott}
  \uses{def:kl-priori-assumptions-canonical-neighborhood-assumption}
  \dcref{Remark 73.8}
  \group{Kleiner--Lott Singularities and Surgery}
 \label{30}
The condition $C_1\geq30\epsilon^{-1}$ ensures that the surgery cap is a
canonical neighborhood.
\end{remark}

\subsection{Evolution of a surgery cap}
\dcref{II.4.5}
\label{secII.4.5}

\begin{lemma}[existence for standard solutions]
  \label{lem:kl-evolution-surgery-cap-existence-standard-solutions}
  \source{kleiner-lott}
  \uses{lem:kl-necks-horns-existence-ricci-flow-surgery}
  \dcref{Lemma II.4.5}
  \group{Kleiner--Lott Singularities and Surgery}
\label{II.4.5}\\
For any $A<\infty$, $\th\in(0,1)$, and $\hat r>0$, there is
$\hat\de=\hat\de(A,\th,\hat r)>0$ such that a Ricci flow with
$(r,\de)$-cutoff on $[a,b]$, with $r(b)\geq\hat r$, and a surgery at
$T_0\in(a,b)$ with $\de(T_0)\leq\hat\de$, has the following property.
If $(p,T_0)\in\M_{T_0}^+$ is the cap center,
$\hat h=h(\de(T_0),\eps,r(T_0),\Phi)$, and
$T_1=\min(b,T_0+\th\hat h^2)$, then either the solution is unscathed on
$P(p,T_0,A\hat h,T_1-T_0)$ and is, after rescaling, $A^{-1}$-close to a
pointed standard solution on
$U_0\times[0,(T_1-T_0)\hat h^{-2}]$, or this holds up to a surgery time
$t^+\in[T_0,T_1)$ and
\[
P(p,T_0,A\hat h,t^+-T_0)\cap\M_{t^+}
\subset\M_{t^+}^--\M_{t^+}^+.
\]
\end{lemma}
\begin{proof}
  \uses{def:kl-priori-assumptions-canonical-neighborhood-assumption,lem:kl-blow-up-time-standard-solution-standard-solutions-lemma,lem:kl-evolution-surgery-cap-existence-standard-solutions-2,lem:kl-length-distortion-estimates-distance-distortion-lemma,lem:kl-necks-horns-existence-ricci-flow-surgery,lem:kl-standard-solutions-existence-kappa-solutions,lem:kl-surgery-pinching-condition-existence-standard-solutions}
\sketch Rescale at the surgery cap. Pointed compactness gives a standard
solution on the maximal interval before the next surgery. If the interval
ended early, bounded curvature and length distortion would keep the cap ball
away from the increasingly long surgery neck, contradicting the occurrence
of another surgery unless the whole ball became extinct.
\end{proof}

\begin{lemma}[existence for standard solutions]
  \label{lem:kl-evolution-surgery-cap-existence-standard-solutions-2}
  \source{kleiner-lott}
  \uses{lem:kl-necks-horns-existence-ricci-flow-surgery}
  \dcref{Lemma 74.2}
  \group{Kleiner--Lott Singularities and Surgery}
 \label{intermediate}
Under the hypotheses of Lemma \ref{II.4.5}, if the solution is unscathed on
$P(p,T_0,A\hat h,T_1-T_0)$, then the pointed region is, after parabolic
rescaling by $\hat h$, $A^{-1}$-close to a pointed subset of a standard
solution on $U_0\times[0,(T_1-T_0)\hat h^{-2}]$.
\end{lemma}
\begin{proof}
  \uses{def:kl-priori-assumptions-canonical-neighborhood-assumption,lem:kl-blow-up-time-standard-solution-standard-solutions-lemma,lem:kl-length-distortion-estimates-distance-distortion-lemma,lem:kl-standard-solutions-existence-kappa-solutions,lem:kl-surgery-pinching-condition-existence-standard-solutions}
\sketch Assume a counterexample sequence with surgery parameters tending to
zero. Rescaled pointed compactness gives a standard-solution limit on a
maximal interval; curvature and length distortion extend the convergence to
the full interval $[0,T_2]$, contradicting the assumed failure of closeness.
\end{proof}

\begin{lemma}[Ricci flow with surgery lemma]
  \label{lem:kl-evolution-surgery-cap-ricci-flow-surgery-lemma}
  \source{kleiner-lott}
  \dcref{Sublemma 3}
  \group{Kleiner--Lott Singularities and Surgery}
$T_3=T_2$.
\end{lemma}
\begin{proof}
  \uses{def:kl-priori-assumptions-canonical-neighborhood-assumption,lem:kl-blow-up-time-standard-solution-standard-solutions-lemma,lem:kl-length-distortion-estimates-distance-distortion-lemma,lem:kl-standard-solutions-existence-kappa-solutions}
\sketch If $T_3<T_2$, a limit on $[0,T_3)$ is a standard solution with
uniform curvature bounds. Canonical-neighborhood and length-distortion
estimates prevent a fixed cap ball from becoming scathed by a sufficiently
long surgery neck, contradicting the definition of $T_3$.
\end{proof}
\begin{proof}
  \proves{lem:kl-evolution-surgery-cap-existence-standard-solutions-2}
\sketch Apply the preceding equality to the counterexample sequence; the
rescaled regions converge to a standard solution through time $T_2$, which
contradicts failure of the asserted closeness.
\end{proof}
\begin{proof}
  \proves{lem:kl-evolution-surgery-cap-existence-standard-solutions}
\sketch Apply the intermediate lemma until the first later surgery. If the
parabolic region is scathed, the surgery neck cannot meet its terminal slice
when $\hat\delta\ll A^{-1}$; hence the entire cap ball is discarded.
\end{proof}

\subsection{Curves that penetrate the surgery region}
\dcref{II.4.6}
\label{II.4.6}

\begin{corollary}[Ricci flow with surgery corollary]
  \label{cor:kl-curves-that-penetrate-surgery-region-ricci-flow-surgery-corollary}
  \source{kleiner-lott}
  \uses{lem:kl-evolution-surgery-cap-existence-standard-solutions, lem:kl-necks-horns-existence-ricci-flow-surgery}
  \dcref{Corollary II.4.6}
  \group{Kleiner--Lott Singularities and Surgery}
\label{corollaryII.4.6}
For every $l<\infty$ and $\hat r>0$ there are $A=A(l,\hat r)<\infty$ and
$\th=\th(l,\hat r)$ such that, in the situation of Lemma \ref{II.4.5},
if $\delta(T_0)<\hat\delta(A,\th,\hat r)$ and
$\ga:[T_0,T_\ga]\ra\M$ is admissible with
$\ga(T_0)\in B(p,T_0,A\hat h/2)$,
$\ga([T_0,T_\ga))\subset P(p,T_0,A\hat h,T_\ga-T_0)$, and either
\begin{enumerate}
\item $T_\ga=T_1=T_0+\th\hat h^2$, or
\item $\ga(T_\ga)\in\D B(p,T_0,A\hat h)\times[T_0,T_\ga]$,
\end{enumerate}
then
\begin{equation}
\int_{T_0}^{T_\ga}\left(R(\gamma(t),t)+|\dot\ga(t)|^2\right)\,dt>l.
\end{equation}
\end{corollary}
\begin{proof}
  \uses{lem:kl-canonical-neighborhood-property-standard-solutions-existence-canonical-neighborhoods,lem:kl-evolution-surgery-cap-existence-standard-solutions}
\sketch Rescale the parabolic region to a standard solution. If the endpoint
time is $\th$, the standard-solution scalar-curvature lower bound makes the
integral diverge as $\th\ra1$. Otherwise the endpoint leaves the rescaled
ball, and nonnegative Ricci curvature plus Cauchy--Schwarz makes the kinetic
term exceed $l$ for $A$ large.
\end{proof}

\subsection{A technical estimate}
\dcref{II.4.7}

\begin{corollary}[existence for a technical estimate]
  \label{cor:kl-technical-estimate-existence-technical-estimate}
  \source{kleiner-lott}
  \uses{lem:kl-evolution-surgery-cap-existence-standard-solutions}
  \dcref{Equation 76.2}
  \group{Kleiner--Lott Singularities and Surgery}
For every $Q<\infty$ and $\hat r>0$, there is
$\th=\th(Q,\hat r)\in(0,1)$ such that, in the situation of Lemma
\ref{II.4.5} with $\delta(T_0)<\hat\delta(A,\th,\hat r)$ and
$A>\eps^{-1}$, a static curve $\ga:[T_0,T_x]\ra\M$ starting in
$B(p,T_0,A\hat h)$ and satisfying
\begin{equation}
\label{Rbound}
Q^{-1}R(\ga(t))\leq R(\ga(T_x))\leq Q(T_x-T_0)^{-1}
\end{equation}
for $t\in[T_0,T_x]$ must have $T_x\leq T_0+\th\hat h^2$.
\end{corollary}

\begin{remark}[estimates for standard solutions]
  \label{rem:kl-technical-estimate-estimates-standard-solutions}
  \source{kleiner-lott}
  \uses{cor:kl-technical-estimate-existence-technical-estimate}
  \dcref{Remark 2}
  \group{Kleiner--Lott Singularities and Surgery}
The hypothesis (\ref{Rbound}) bounds the scalar curvature along the static
curve in the endpoint curvature scale and bounds its elapsed time by the
corresponding rescaled standard-solution lifetime.
\end{remark}
\begin{proof}
  \uses{cor:kl-technical-estimate-existence-technical-estimate,lem:kl-canonical-neighborhood-property-standard-solutions-existence-canonical-neighborhoods,lem:kl-evolution-surgery-cap-existence-standard-solutions}
\sketch If $T_x>T_0+\th\hat h^2$, the standard-solution lower bound gives
\[
R(\ga(T_0+\th\hat h^2))\geq\const(1-\th)^{-1}\hat h^{-2}.
\]
Together with (\ref{Rbound}) this yields
\[
T_x-T_0\leq\const Q^2(1-\th)\hat h^2,
\]
which contradicts $T_x-T_0> \th\hat h^2$ when $\th$ is sufficiently close
to $1$.
\end{proof}

\section{Statement of the existence theorem for Ricci flow with surgery}
\dcref{II.5} \label{II.5}

\begin{definition}[normalized definition]
  \label{def:kl-statement-existence-theorem-ricci-flow-surgery-normalized-definition}
  \source{kleiner-lott}
  \dcref{Definition 1}
  \group{Kleiner--Lott Surgery Existence}

A compact Riemannian $3$-manifold is {\em normalized} if
$|\Rm|\leq 1$ everywhere, and the volume of every unit
ball is at least half the volume of the Euclidean unit
ball.
\end{definition}

\begin{proposition}[monotonicity for kappa solutions]
  \label{prop:kl-statement-existence-theorem-ricci-flow-surgery-monotonicity-kappa-solutions}
  \source{kleiner-lott}
  \dcref{Proposition II.5.1}
  \group{Kleiner--Lott Surgery Existence}
(cf. Proposition II.5.1)
\label{surgeryflowexists}
There exist decreasing sequences $0<r_j<\eps^2$,
$\kappa_j>0$, $0<\bar\de_j<\eps^2$ for $1\leq j<\infty$,
such that for any normalized initial data and any nonincreasing function
$\de:[0,\infty)\ra(0,\infty)$ such that $\de<\bar\de_j$
on $[2^{j-1}\eps,2^j\eps]$, the Ricci flow with $(r,\de)$-cutoff
is defined for all time and is $\kappa$-noncollapsed
at scales below $\eps$.
\end{proposition}

\subsection{The $L$-function and Ricci flows with surgery}
\dcref{II.5.1}
\label{lfunctionandsurgerysection}

Let $\M$ be a Ricci flow with surgery,
and fix a point $(x_0,t_0)\in\M$.   One may
define the $\L$-length of an  admissible curve $\ga$ from
$(x_0,t_0)$ to some $(x,t)$, for $t<t_0$, using the formula

\begin{equation}
\L(\ga)=\int_t^{t_0}\sqrt{t_0-\bar t}\,\left(R+|\dot{\ga}|^2\right)d\bar t,
\end{equation}

\begin{equation} \label{newbarrier}
\bar L_\tau+\De \bar L\leq 6
\end{equation}

\begin{lemma}[existence for Ricci flow with surgery]
  \label{lem:kl-reduced-l-geometry-existence-ricci-flow-surgery}
  \source{kleiner-lott}
  \dcref{Lemma 78.3}
  \group{Kleiner--Lott Surgery Existence}
[Existence of $\L$-minimizers]
\label{minimizersexist}
Let $\M$ be a Ricci flow with surgery
defined on $ [a,b]$. Suppose that $(x_0,t_0)\in \M$ lies in the backward
time slice $\M_{t_0}^-$.

(1) For each $(x,t)\in \M_{[a,t_0)}$ with $L(x,t)<\infty$, there exists
an $\L$-minimizing admissible path $\ga:[t,t_0]\ra\M$ from $(x,t)$ to $(x_0,t_0)$
which satisfies the
$\L$-geodesic equation at every time $\overline{t}\in (t,t_0)$ for which
$\ga(\overline{t})\in\M_{\reg}$.


(2) $L$ is lower semicontinuous on
$\M_{[a,t_0)}$
and continuous
on
$\M_{\reg}\cap \M_{[a,t_0)}$. (Note that $\M_a^+ \subset \M_{\reg}$.)

(3) Every sequence $(x_j,t_j)\in \M_{[a,t_0)}$ with $\limsup_j L(x_j,t_j)<\infty$
has  a convergent subsequence.
\end{lemma}

\begin{lemma}[existence for reduced length]
  \label{lem:kl-reduced-l-geometry-existence-reduced-length}
  \source{kleiner-lott}
  \dcref{Lemma 78.6}
  \group{Kleiner--Lott Surgery Existence}

\label{stillleq3/2}
Suppose that $\M$ is a Ricci flow with surgery defined on $ [a,b]$.
Take
$t_0\in (a,b]$
and $(x_0,t_0)\in \M_{t_0}^-$.  Suppose that
for every $t \in [a, t_0)$, every admissible curve
$[t,t_0]\ra\M$ ending at $(x_0,t_0)$ which does not lie
in
$\M_{\reg}\cup  \M_{t_0}^-$
has reduced length strictly greater than $\frac{3}{2}$. Then
there is a point $(x,a)\in \M_a^+$ where $l(x,a)\leq \frac{3}{2}$.
\end{lemma}

\begin{remark}[structural remark on reduced length]
\label{rem:kl-reduced-l-geometry-structural-remark-reduced-length}
\source{kleiner-lott}
\dcref{Remark 4}
\group{Kleiner--Lott Surgery Existence}


The hypothesis concerns reduced lengths of barely admissible curves, including curves that meet a surgery slice.
\end{remark}

\begin{proof}
\sketch Use the reduced-length and reduced-volume identities, the maximum principle, and the cited surgery estimates; the resulting barriers give the stated bound.
\uses{lem:kl-reduced-l-geometry-existence-ricci-flow-surgery,lem:kl-reduced-l-geometry-positivity-reduced-l-geometry}
\end{proof}

\begin{definition}[kappa noncollapsing]
  \label{def:kl-reduced-l-geometry-kappa-noncollapsing}
  \source{kleiner-lott}
  \dcref{Definition 78.10}
  \group{Kleiner--Lott Surgery Existence}

\label{surgerycollapsedef}
Let $\M$ be a Ricci flow with surgery defined on $[a,b]$. Suppose that
$(x_0,t_0)\in\M$ and $r>0$ are such that $t_0-r^2\geq a$,
$B(x_0,t_0,r)\subset\M_{t_0}^-$ is a proper ball
and the parabolic ball $P(x_0,t_0,r,-r^2)$ is unscathed.
Then {\em $\M$ is $\kappa$-collapsed at $(x_0,t_0)$ at scale
$r$} if $|\Rm|\leq r^{-2}$ on $P(x_0,t_0,r,-r^2)$ and
$\vol(B(x_0,t_0,r))<\kappa r^3$; otherwise it is $\kappa$-noncollapsed.
\end{definition}

\begin{lemma}[existence for reduced volume]
  \label{lem:kl-reduced-l-geometry-existence-reduced-volume}
  \source{kleiner-lott}
  \dcref{Lemma 78.11}
  \group{Kleiner--Lott Surgery Existence}
(Local version of reduced volume comparison)
\label{localreducedvolume}
There is a function $\kappa':\R_+\ra \R_+$, satisfying
$\lim_{\kappa\ra 0}\kappa'(\kappa)=0$, with the following property.
Let $\M$ be a Ricci flow with surgery defined on $[a,b]$. Suppose
that we are given $t_0\in (a,b]$,
$(x_0,t_0)\in \M_{t_0}\cap \M_{\reg}$, $t\in [a,t_0)$ and
$r \in (0,\sqrt{t_0-t})$.  Let $Y$ be the set of points $(x,t)\in \M_t$
that are accessible from $(x_0,t_0)$ by means of minimizing $\L$-geodesics
which remain in $\M_{\reg}$.  Assume
in addition that
$\M$ is $\kappa$-collapsed at $(x_0,t_0)$ at scale $r$, i.e.
$P(x_0,t_0,r,-r^2) \cap \M_{[t_0 - r^2, t_0)} \subset\M_{\reg}$,
$|\Rm|\leq r^{-2}$ on
$P(x_0,t_0,r,-r^2)$, and $\vol(B(x_0,t_0,r))<\kappa r^3$.
Then  the reduced volume of $Y$ is at most
$\kappa'(\kappa)$.
\end{lemma}

\begin{proof}
\sketch Use the reduced-length and reduced-volume identities, the maximum principle, and the cited surgery estimates; the resulting barriers give the stated bound.
\uses{thm:kl-no-local-collapsing-theorem-ii-existence-kappa-solutions}
\end{proof}

\subsection{Establishing noncollapsing in the presence of surgery}
\dcref{II.5.2}
\label{II.5.2section}

\begin{definition}[noncollapsing definition]
  \label{def:kl-establishing-noncollapsing-presence-surgery-noncollapsing-definition}
  \source{kleiner-lott}
  \dcref{Definition 1}
  \group{Kleiner--Lott Surgery Existence}

The {\em ${\mathcal L}_+$-length} of an admissible curve $\gamma$ is
\begin{equation}
{\mathcal L}_+(\gamma, \tau) \: = \:
\int_{t_0 - \tau}^{t_0} \sqrt{t_0 - t} \: \left(
R_+(\gamma(t), t) \: + \: |\dot{\gamma}(t)|^2 \right) \: dt,
\end{equation}
where $R_+(x,t) = \max(R(x,t), 0)$.
\end{definition}

\begin{lemma}[existence for noncollapsing]
  \label{lem:kl-establishing-noncollapsing-presence-surgery-existence-noncollapsing}
  \source{kleiner-lott}
  \dcref{Lemma 79.3}
  \group{Kleiner--Lott Surgery Existence}
(Forcing $\L_+$ to be large, cf. Lemma II.5.3)
\label{lplusbig}

For all $\La< \infty$, $\bar r>0$ and $\hat r>0$,
there is  a constant  $F_0=F_0(\La,\bar r,\hat r)$
with the following property.  Suppose that

$\bullet$  $\M$ is a Ricci flow  with $(r,\de)$-cutoff
defined on an interval containing $[t,t_0]$, where
 $r([t,t_0])\subset [\hat r,\eps]$,

$\bullet$ $r_0\geq \bar r$,
 $B(x_0,t_0,r_0)$ is a proper ball which is unscathed on $[t_0-r_0^2,t_0]$, and
$|\Rm|\leq r_0^{-2}$ on $P(x_0,t_0,r_0,-r_0^2)$,

$\bullet$ $\ga:[t,t_0]\ra\M$ is an admissible curve
ending at $(x_0, t_0)$
whose image
is not contained in $\M_{\reg}\cup\M_{t_0}$, and

$\bullet$ $\de < F_0(\La,\bar r,\hat r)$
on $[t,t_0]$.

\noindent
Then $\L_+(\ga)>\La$.
\end{lemma}

\begin{proof}
\sketch Use the reduced-length and reduced-volume identities, the maximum principle, and the cited surgery estimates; the resulting barriers give the stated bound.
\uses{cor:kl-curves-that-penetrate-surgery-region-ricci-flow-surgery-corollary,lem:kl-evolution-surgery-cap-existence-standard-solutions}

\begin{equation}
\label{h<<r}
\max_{[t,t_0]}h(t)\leq \left(\max_{[t,t_0]}\de\right)^2
\left(\max_{[t,t_0]}r(t)\right)\leq F_0^2\eps,
\end{equation}

\begin{equation}
\label{aandthetaconditions}
A=A((\De t)^{-\frac12}\La,\hat r),\quad
\th=\th((\De t)^{-\frac12}\La,\hat r)
\end{equation}

\begin{equation}
\label{deltasmall}
 F_0\leq \hat\de(A+2,\th,\hat r)
\end{equation}

\begin{equation}
\label{p'sdisjoint}
P(x_0,t_0,r_0,-r_0^2)\cap P(x,t,Ah(t),\th h^2(t))=\emptyset
\end{equation}
\end{proof}

\begin{lemma}[vanishing for noncollapsing]
  \label{lem:kl-establishing-noncollapsing-presence-surgery-vanishing-noncollapsing}
  \source{kleiner-lott}
  \dcref{Lemma 79.11}
  \group{Kleiner--Lott Surgery Existence}
 \label{Rev}
If $\M$ is a Ricci flow with surgery, with normalized initial
condition at time zero, then for all $t \ge 0$, $R(x,t) \: \ge \:
- \: \frac{3}{2} \: \frac{1}{t+\frac14}$.
\end{lemma}

\begin{proof}
\sketch Use the reduced-length and reduced-volume identities, the maximum principle, and the cited surgery estimates; the resulting barriers give the stated bound.
\end{proof}

\begin{lemma}[existence for kappa solutions]
  \label{lem:kl-establishing-noncollapsing-presence-surgery-existence-kappa-solutions}
  \source{kleiner-lott}
  \dcref{Lemma II.5.2}
  \group{Kleiner--Lott Surgery Existence}
(Noncollapsing estimate) (cf. Lemma II.5.2)
\label{II.5.2}

Suppose $\epsilon \ge r_- \ge r_+ >0$, $\kappa_->0$, $E_->0$ and $E<\infty$.
Then there are constants $\overline{\delta} =
\overline{\delta}(r_-,r_+,\kappa_-,E_-,E)$
and $\kappa_+=\kappa_+(r_-,\kappa_-,E_-,E)$ with
the following property.  Suppose that

$\bullet$  $a<b<c,\quad b-a\geq E_-,\quad c-a\leq E$,

$\bullet$  $\M$ is a Ricci flow with $(r,\de)$-cutoff
with normalized initial condition
defined on a time interval containing $[a,c)$,

$\bullet$ $r\geq r_-$ on $[a,b)$ and $r\geq r_+$ on $[b,c)$,

$\bullet$ $r \le \epsilon$ ,

$\bullet$ $\M$ is $\kappa_-$-noncollapsed at scales below
$\eps$ on $[a,b)$ and

$\bullet$ $\de\leq \bar\de$ on $[a,c)$,


Then $\M$ is $\kappa_+$-noncollapsed at scales below $\eps$
on $[b,c)$.
\end{lemma}

\begin{remark}[estimates for kappa solutions]
\label{rem:kl-establishing-noncollapsing-presence-surgery-estimates-kappa-solutions}
\source{kleiner-lott}
\dcref{Remark 5}
\group{Kleiner--Lott Surgery Existence}


The parameter $\bar\delta$ may depend on the lower bound $r_+$, while the resulting noncollapsing constant $\kappa_+$ is independent of $r_+$.
\end{remark}

\begin{proof}
\sketch Argue by contradiction and rescale at the first failure. Hamilton--Ivey pinching, noncollapsing, and compactness yield a $\kappa$-solution or standard-solution limit, contradicting the failure.
\uses{def:kl-priori-assumptions-canonical-neighborhood-assumption,lem:kl-canonical-neighborhood-property-standard-solutions-existence-canonical-neighborhoods,lem:kl-curvature-bounds-from-priori-assumptions-canonical-neighborhoods-lemma,lem:kl-curvature-bounds-from-priori-assumptions-estimates-curvature,lem:kl-establishing-noncollapsing-presence-surgery-existence-kappa-solutions,lem:kl-establishing-noncollapsing-presence-surgery-existence-noncollapsing,lem:kl-establishing-noncollapsing-presence-surgery-vanishing-noncollapsing,lem:kl-evolution-surgery-cap-existence-standard-solutions,lem:kl-performing-surgery-continuing-flows-existence-canonical-neighborhoods,lem:kl-reduced-l-geometry-existence-reduced-length,lem:kl-reduced-l-geometry-existence-reduced-volume,lem:kl-reduced-l-geometry-existence-ricci-flow-surgery,lem:kl-standard-solutions-existence-kappa-solutions,lem:kl-structure-at-first-singularity-time-characterization-structure-at-first-singularity-time,prop:kl-statement-existence-theorem-ricci-flow-surgery-monotonicity-kappa-solutions,thm:kl-canonical-neighborhood-theorem-compactness-canonical-neighborhoods,thm:kl-no-local-collapsing-theorem-ii-existence-kappa-solutions}

\begin{gather}
\label{llplus}
\L(\ga)\leq \L_+(\ga)\leq \L(\ga)+\int_a^c6\sqrt{c-t}\: dt
\leq \L(\ga)+4E^{\frac32} \\
\mbox{and}\quad l(x,t)\geq\frac{L_+-4E^{\frac32}}{2E^{\frac12}}. \notag
\end{gather}

\begin{equation}
\label{lplusupper}
\L_+(\ga)\leq 3\sqrt{t_0-a}+4E^{\frac32}\leq 3\sqrt{E}+4E^{\frac32}.
\end{equation}

\begin{equation}
\label{Rbound2}
R\leq \const s^{-2}
\end{equation}

\begin{equation}
\label{kappa1}
\vol(B(x_0,t_0,r_0))r_0^{-3}\geq \kappa_1=\kappa_1(r_-,\kappa_-,E_-,E).
\end{equation}
\end{proof}

\begin{lemma}[kappa solutions lemma]
  \label{lem:kl-establishing-noncollapsing-presence-surgery-kappa-solutions-lemma}
  \source{kleiner-lott}
  \dcref{Lemma 79.23}
  \group{Kleiner--Lott Surgery Existence}

\label{r0/100}
If $r_0< \frac{r_+}{100}$ then $\vol(B(x_0,t_0,r_0))r_0^{-3}\geq
\kappa_3=\kappa_3(r_-,\kappa_-,E_-,E)$.
\end{lemma}

\begin{proof}
\sketch Argue by contradiction and rescale at the first failure. Hamilton--Ivey pinching, noncollapsing, and compactness yield a $\kappa$-solution or standard-solution limit, contradicting the failure.
\uses{def:kl-priori-assumptions-canonical-neighborhood-assumption,lem:kl-establishing-noncollapsing-presence-surgery-existence-kappa-solutions,lem:kl-standard-solutions-existence-kappa-solutions}
\proves{rem:kl-establishing-noncollapsing-presence-surgery-estimates-kappa-solutions}
The retained proof uses \cite[Corollary 7.7]{Chow-Knopf}.

\begin{equation}
\label{F_1condition}
\bar\de^{\al}<\bar\de(r_i,r^\al,\kappa_i,2^{i-1}\eps,3(2^{i-1})\eps)
\end{equation}
\end{proof}

\subsection{Construction of the Ricci flow with surgery}  \label{surgeryflow}

\begin{lemma}[Ricci flow with surgery lemma]
  \label{lem:kl-construction-ricci-flow-surgery-ricci-flow-surgery-lemma}
  \source{kleiner-lott}
  \dcref{Sublemma 1}
  \group{Kleiner--Lott Surgery Existence}

For all $\la<\infty$,
the ball $B(\bar x^\al,0,\la)\subset\widehat\M^\al_0$
is proper for sufficiently
large $\al$.
\end{lemma}

\begin{proof}
\sketch Apply the preceding estimates and compactness argument to the stated hypotheses; the limiting or inductive construction has the required properties.
\uses{lem:kl-curvature-bounds-from-priori-assumptions-canonical-neighborhoods-lemma,lem:kl-structure-at-first-singularity-time-characterization-structure-at-first-singularity-time}
\proves{rem:kl-establishing-noncollapsing-presence-surgery-estimates-kappa-solutions}
\end{proof}

\begin{lemma}[convergence for Ricci flow with surgery]
  \label{lem:kl-construction-ricci-flow-surgery-convergence-ricci-flow-surgery}
  \source{kleiner-lott}
  \dcref{Lemma 80.4}
  \group{Kleiner--Lott Surgery Existence}
 \label{bdrysmoothness}
After passing to a subsequence
if necessary, the pointed flows  $(\widehat\M^\al,(\bar x^\al,0))$
converge on the time interval $(T_1, 0]$ to a Ricci flow (without
surgery) $(\M^\infty,(\bar x^\infty,0))$ with
a smooth complete nonnegatively-curved Riemannian metric
on each time slice, and
scalar curvature globally bounded above by
some number $Q<\infty$.
(We interpret $(0, 0]$ to mean $\{0\}$ rather than the
empty set.)
\end{lemma}

\begin{proof}
\sketch The curvature and noncollapsing bounds give pointed compactness; Hamilton--Ivey pinching and the surgery estimates identify the smooth limit and preserve the stated properties.
\uses{lem:kl-curvature-bounds-from-priori-assumptions-canonical-neighborhoods-lemma,lem:kl-curvature-bounds-from-priori-assumptions-estimates-curvature,lem:kl-evolution-surgery-cap-existence-standard-solutions,prop:kl-statement-existence-theorem-ricci-flow-surgery-monotonicity-kappa-solutions,thm:kl-canonical-neighborhood-theorem-compactness-canonical-neighborhoods}
\proves{rem:kl-establishing-noncollapsing-presence-surgery-estimates-kappa-solutions}
The retained proof uses \cite{Perelman2}.

\begin{equation}
\label{2q+1}
R(x,t)< 8(Q+2)\quad\mbox{for all}\quad (x,t)\in P(\bar x^\al,0,\la,t^\al).
\end{equation}
\end{proof}

\section{Double sided curvature bound in the thick part}
\dcref{II.6}
\label{II.6}

\begin{lemma}[compactness for thick--thin geometry]
  \label{lem:kl-double-sided-curvature-bound-thick-part-compactness-thick-thin-geometry}
  \source{kleiner-lott}
  \dcref{Lemma 81.1}
  \group{Kleiner--Lott Long-Time Curvature}
 \label{pscempty}
If $\M$ is a Ricci flow with $(r,\de)$-cutoff on a compact manifold and
$g(0)$ has positive scalar curvature then the solution goes extinct
after a finite time, i.e. $\M_T \: = \: \emptyset$ for some $T > 0$.
\end{lemma}

\begin{proof}
\sketch The curvature and noncollapsing bounds give pointed compactness; Hamilton--Ivey pinching and the surgery estimates identify the smooth limit and preserve the stated properties.
\uses{lem:kl-establishing-noncollapsing-presence-surgery-vanishing-noncollapsing}
\end{proof}

\begin{lemma}[compactness for Ricci flow with surgery]
  \label{lem:kl-double-sided-curvature-bound-thick-part-compactness-ricci-flow-surgery}
  \source{kleiner-lott}
  \dcref{Lemma 81.2}
  \group{Kleiner--Lott Long-Time Curvature}
 \label{finitetime}
If $\M$ is a Ricci flow with surgery that goes extinct after a finite time,
then the initial (compact connected orientable) $3$-manifold is diffeomorphic to a
connected sum of $S^1 \times S^2$'s and quotients of the round
$S^3$.
\end{lemma}

\begin{proof}
\sketch The curvature and noncollapsing bounds give pointed compactness; Hamilton--Ivey pinching and the surgery estimates identify the smooth limit and preserve the stated properties.
\uses{lem:kl-performing-surgery-continuing-flows-rigidity-ricci-flow-surgery}
\end{proof}

\begin{corollary}[existence for Ricci flow with surgery]
  \label{cor:kl-double-sided-curvature-bound-thick-part-existence-ricci-flow-surgery}
  \source{kleiner-lott}
  \dcref{Corollary II.6.8}
  \group{Kleiner--Lott Long-Time Curvature}
 \label{II.6.8}
(cf. Corollary II.6.8) For any $w > 0$ one can find
$\tau = \tau(w) > 0$, $K = K(w) < \infty$,
$\overline{r} = \overline{r}(w) > 0$ and $\theta = \theta(w) > 0$ with
the following property. Suppose we have a solution to the Ricci flow
with $(r, \delta)$-cutoff on the time interval $[0, t_0]$, with
normalized initial data.
Let $h_{\max}(t_0)$ be the maximal surgery radius on $[t_0/2, t_0]$.
(If there are no surgeries on $[t_0/2, t_0]$ then $h_{\max}(t_0) = 0$.)
Let $r_0$ satisfy \\
1. $\theta^{-1}(w) h_{\max}(t_0) \: \le \: r_0 \: \le \:
\overline{r} \sqrt{t_0}$.\\
2. The ball $B(x_0, t_0, r_0)$ has sectional curvatures at least
$- \: r_0^{-2}$ at each point.\\
3. $\vol(B(x_0, t_0, r_0)) \: \ge \: w r_0^3$.\\

Then the solution is unscathed in $P(x_0, t_0, r_0/4, - \tau r_0^2)$
and satisfies $R \: < \: Kr_0^{-2}$ there.
\end{corollary}

A more substantial difference is that the smoothness of the flow
in Corollary \ref{corI.12.4}
is guaranteed by the setup, whereas in Corollary \ref{II.6.8} we must
prove that the solution is unscathed in $P(x_0, t_0, r_0/4, - \tau r_0^2)$.

\subsection{Earlier scalar curvature bounds on smaller balls from lower curvature bounds and a later volume bound}
\dcref{II.6.5} \label{II.6.5}

\begin{enumerate}
\item If $t \notin [c,d]$ then ${\mathcal U} \cap \M_t = \emptyset$.
\item If $t \in [c,d]$ is not
a singularity time then ${\mathcal U} \cap \M_t$ is the $r$-ball around $x$ in
$\M_t$.
\item If $t \in [c,d]$ is a surgery time $t_k^+$ then
${\mathcal U} \cap \M_t$ is
the image in $\M_t$ of an $r$-ball around $x$ in
$\Omega_k$, that lies entirely in $X_k^+ \subset \Omega_k$.
\end{enumerate}

\begin{lemma}[estimates for scalar curvature]
  \label{lem:kl-earlier-scalar-curvature-bounds-smaller-balls-from-lower-curvature-bound-estimates-scalar-curvature}
  \source{kleiner-lott}
  \dcref{Lemma 82.1}
  \group{Kleiner--Lott Long-Time Curvature}
 \label{I.11.6(a)}
Given $w > 0$, there exist $\tau_0 = \tau_0(w) > 0$ and
$K_0^\prime = K_0^\prime(w) < \infty$
with the following property. Suppose that we have
a Ricci flow with $(r,\de)$-cutoff such that \\
1. $\bigcup_{t \in
[-\tau r_0^2, 0]} B(x_0, t, r_0)$ is an unscathed time-dependent family of
metric balls, where $\tau \: \le \: \tau_0$. \\
2. The sectional curvatures are bounded below by $- \: r_0^{-2}$
on the
above family of balls. \\
3. $\vol(B(x_0, 0, r_0)) \: \ge \: w r_0^3$. \\

Then
\begin{enumerate}
\item
$R \: \le \: K_0^\prime \: \tau^{-1} \: r_0^{-2}$
on $\bigcup_{ t \in [- \tau r_0^2/2, 0]} B(x_0, t, r_0/2)$, and
\item $\vol(B(x_0, -\tau r_0^2, r_0/2))$ is at least $\frac{1}{10}$
of the volume of the Euclidean ball of the same radius.
\end{enumerate}
\end{lemma}

\begin{lemma}[estimates for scalar curvature]
  \label{lem:kl-earlier-scalar-curvature-bounds-smaller-balls-from-lower-curvature-bound-estimates-scalar-curvature-4}
  \source{kleiner-lott}
  \dcref{Lemma II.6.5(a)}
  \group{Kleiner--Lott Long-Time Curvature}
 \label{II.6.5(a)}
(cf. Lemma II.6.5(a))
Given $w > 0$, there exist $\tau_0 = \tau_0(w) > 0$ and $K_0 =
K_0(w) < \infty$ with the following property. Suppose that we have
a Ricci flow with $(r,\de)$-cutoff such that \\
1. The parabolic neighborhood
$P(x_0, 0, r_0, -\tau r_0^2)$ is unscathed,
where $\tau \: \le \: \tau_0$. \\
2. The sectional curvatures are bounded below by $- \: r_0^{-2}$ on
$P(x_0, 0, r_0, -\tau r_0^2)$. \\
3. $\vol(B(x_0, 0, r_0)) \: \ge \: wr_0^3$. \\

Then $R \: \le \: K_0 \: \tau^{-1} \: r_0^{-2}$ on
$P(x_0, 0, r_0/4, -\tau r_0^2/2)$.
\end{lemma}

\begin{proof}
\sketch Use the lower curvature and volume hypotheses with point selection, distance distortion, and local curvature estimates; a contradiction argument gives the stated upper bound.
\uses{lem:kl-earlier-scalar-curvature-bounds-smaller-balls-from-lower-curvature-bound-estimates-scalar-curvature,lem:kl-length-distortion-estimates-distance-distortion-lemma}
\end{proof}

There is an evident analogy between Lemma \ref{II.6.5(a)} and
Corollary \ref{II.6.8}. However, there is the important difference that
Corollary \ref{II.6.8} (along with Corollary \ref{corI.12.4}) only
assumes a lower sectional curvature bound at the final time slice.

\subsection{Locating small balls whose subballs have almost Euclidean volume}
\dcref{II.6.6}

\begin{lemma}[existence for sectional curvature]
  \label{lem:kl-locating-small-balls-whose-subballs-almost-euclidean-volume-existence-sectional-curvature}
  \source{kleiner-lott}
  \dcref{Lemma II.6.6}
  \group{Kleiner--Lott Long-Time Curvature}
 \label{Alexlemma2} (cf. Lemma II.6.6)
For any $\widehat{\epsilon}, w>0$
there exists $\theta_0 = \theta_0(\widehat{\epsilon},w)$
such that if $B(x,1)$ is a metric ball of volume at least $w$, compactly
contained in a manifold without boundary with sectional
curvatures at least $-1$, then there exists a subball
$B(y,\theta_0) \subset B(x,1)$ such that every subball
$B(z,r) \subset B(y,\theta_0)$ of any radius has volume at
least $(1-\widehat{\epsilon})$ times the volume of the Euclidean ball
of the same radius.
\end{lemma}

\begin{remark}[existence for volume]
\label{rem:kl-locating-small-balls-whose-subballs-almost-euclidean-volume-existence-volume}
\source{kleiner-lott}
\uses{lem:kl-locating-small-balls-whose-subballs-almost-euclidean-volume-existence-sectional-curvature}
\dcref{Remark 2}
\group{Kleiner--Lott Long-Time Curvature}


\end{remark}

\subsection{Proof of the double sided curvature bound in the thick part, modulo two propositions}
\dcref{II.6.8}

\begin{proposition}[existence for canonical neighborhoods]
  \label{prop:kl-proof-double-sided-curvature-bound-thick-part-modulo-two-propositions-existence-canonical-neighborhoods}
  \source{kleiner-lott}
  \uses{def:kl-priori-assumptions-canonical-neighborhood-assumption}
  \dcref{Proposition II.6.3}
  \group{Kleiner--Lott Long-Time Curvature}
 \label{propII.6.3}
(cf. Proposition II.6.3)
For any
$A < \infty$
one can find positive constants
$\kappa(A)$, $K_1(A)$, $K_2(A)$, $\overline{r}(A)$,
such that for any $t_0 < \infty$
there exists $\overline{\delta}_{A}(t_0)
> 0$, decreasing in $t_0$,
with the following property.  Suppose that
we have a Ricci flow with $(r, \delta)$-cutoff on a time interval
$[0, T]$, where $\delta(t) < \overline{\delta}_A(t)$ on $[t_0/2, t_0]$, with
normalized initial data.  Assume that \\
1. The solution is unscathed on a parabolic ball
$P(x_0, t_0, r_0, -r_0^2)$, with $2r_0^2 < t_0$. \\
2. $|\Rm| \le \frac{1}{3r_0^2}$ on $P(x_0, t_0, r_0, -r_0^2)$. \\
3. $\vol(B(x_0, t_0, r_0)) \: \ge \: A^{-1} r_0^3$.\\

Then\\
(a) The solution is $\kappa$-noncollapsed on scales less than $r_0$
in $B(x_0, t_0, Ar_0)$.\\
(b) Every point $x \in B(x_0, t_0, Ar_0)$ with $R(x, t_0) \ge K_1 r_0^{-2}$
has a canonical neighborhood in the sense of
Definition \ref{canonicalnbhddef}.\\
(c) If $r_0 \: \le \: \overline{r} \sqrt{t_0}$ then $R \: \le \: K_2
r_0^{-2}$ in $B(x_0, t_0, Ar_0)$.
\end{proposition}

\begin{proposition}[existence for thick--thin geometry]
  \label{prop:kl-proof-double-sided-curvature-bound-thick-part-modulo-two-propositions-existence-thick-thin-geometry}
  \source{kleiner-lott}
  \dcref{Proposition II.6.4}
  \group{Kleiner--Lott Long-Time Curvature}
 \label{propII.6.4}
(cf. Proposition II.6.4)
There exist $\tau, \overline{r}, C_1 > 0$ and $K < \infty$ with the
following property. Suppose that we have a Ricci flow with $(r, \delta)$-cutoff
on the time interval $[0, t_0]$, with normalized initial data.
  Let $r_0$ satisfy
$2C_1 h_{\max}(t_0) \: \le \: r_0 \: \le \: \overline{r} \sqrt{t_0}$,
where $h_{\max}(t_0)$ is the maximal cutoff radius for surgeries in
$[t_0/2, t_0]$. (If there are no surgeries on $[t_0/2, t_0]$ then
$h_{\max}(t_0) = 0$.)

Assume \\
1. The ball $B(x_0, t_0, r_0)$ has sectional curvatures at least $- r_0^{-2}$
at each point. \\
2. The volume of any subball $B(x, t_0, r) \subset B(x_0, t_0, r_0)$ with
any radius $r > 0$ is at least $(1-\epsilon)$ times the volume of the
Euclidean ball of the same radius.\\

Then the solution is unscathed on
$P(x_0, t_0, r_0/4, - \tau r_0^2)$
and satisfies $R \: < \: K r_0^{-2}$ there.
\end{proposition}

\subsection{Canonical neighborhoods and later curvature bounds on bigger balls from curvature and volume bounds}
\dcref{II.6.3}

\begin{proposition}[existence for canonical neighborhoods]
  \label{prop:kl-canonical-neighborhoods-later-curvature-bounds-bigger-balls-from-curvatu-existence-canonical-neighborhoods}
  \source{kleiner-lott}
  \uses{def:kl-priori-assumptions-canonical-neighborhood-assumption}
  \dcref{Proposition 1}
  \group{Kleiner--Lott Long-Time Curvature}

(cf. Proposition II.6.3)
For any $A > 0$ one can find positive constants $\kappa(A)$, $K_1(A)$, $K_2(A)$,
$\overline{r}(A)$,
such that for any $t_0 < \infty$
there exists $\overline{\delta}_{A}(t_0)
> 0$, decreasing in $t_0$, with the following property.  Suppose that
we have a Ricci flow with $(r, \delta)$-cutoff on a time interval
$[0, T]$, where $\delta(t) < \overline{\delta}_A(t)$ on $[t_0/2, t_0]$, with
normalized initial data.  Assume that \\
1. The solution is unscathed on a parabolic neighborhood
$P(x_0, t_0, r_0, -r_0^2)$, with $2r_0^2 < t_0$. \\
2.
$|\Rm| \le \frac13 r_0^{-2}$
on $P(x_0, t_0, r_0, -r_0^2)$. \\
3. $\vol(B(x_0, t_0, r_0)) \: \ge \: A^{-1} r_0^3$.\\

Then\\
(a) The solution is $\kappa$-noncollapsed on scales less than $r_0$
in $B(x_0, t_0, Ar_0)$.\\
(b) Every point $x \in B(x_0, t_0, Ar_0)$ with $R(x, t_0) \ge K_1 r_0^{-2}$
has a canonical neighborhood in the sense of
Definition \ref{canonicalnbhddef}.\\
(c) If $r_0 \: \le \: \overline{r} \sqrt{t_0}$ then $R \: \le \: K_2
r_0^{-2}$ in $B(x_0, t_0, Ar_0)$.
\end{proposition}

\subsection{Earlier scalar curvature bounds on smaller balls from lower curvature bounds and volume bounds, in the presence of possible surgeries}
\dcref{II.6.4}

\begin{proposition}[existence for scalar curvature]
  \label{prop:kl-earlier-scalar-curvature-bounds-smaller-balls-from-lower-curvature-bound-existence-scalar-curvature}
  \source{kleiner-lott}
  \dcref{Proposition 5}
  \group{Kleiner--Lott Long-Time Curvature}

There exist $\tau, \overline{r}, C_1 > 0$ and $K < \infty$
with the
following property. Suppose that we have a Ricci flow with $(r, \delta)$-cutoff
on the time interval $[0, t_0]$, with normalized initial data.  Let
$r_0$ satisfy
$2C_1 h_{\max}(t_0) \: \le \: r_0 \: \le \: \overline{r} \sqrt{t_0}$,
where $h_{\max}(t_0)$ is the maximal cutoff radius for surgeries in
$[t_0/2, t_0]$. (If there are no surgeries on $[t_0/2, t_0]$ then we put
$h_{\max}(t_0) = 0$.)
Assume \\
1. The ball $B(x_0, t_0, r_0)$ has sectional curvatures at least $- r_0^{-2}$
at each point. \\
2. The volume of any subball $B(x, t_0, r) \subset B(x_0, t_0, r_0)$ with
any radius $r > 0$ is at least $(1-\epsilon)$ times the volume of the
Euclidean ball of the same radius.\\

Then the solution is unscathed on
$P(x_0, t_0, r_0/4, - \tau r_0^2)$ and satisfies $R < K r_0^{-2}$
there.
\end{proposition}

\begin{proof}
\sketch Use the lower curvature and volume hypotheses with point selection, distance distortion, and local curvature estimates; a contradiction argument gives the stated upper bound.
\end{proof}

\begin{lemma}[existence for scalar curvature]
  \label{lem:kl-earlier-scalar-curvature-bounds-smaller-balls-from-lower-curvature-bound-existence-scalar-curvature-6}
  \source{kleiner-lott}
  \dcref{Lemma 86.2}
  \group{Kleiner--Lott Long-Time Curvature}
 \label{extra}
There exists $\tau^\prime > 0$
with the
following property. Suppose that we have a Ricci flow with $(r, \delta)$-cutoff
on the time interval $[0, t_0]$, with normalized initial data.  Let
$r_0$ satisfy
$2C_1 h_{\max}(t_0) \: \le \: r_0 \: \le \: \overline{r} \sqrt{t_0}$,
where $h_{\max}(t_0)$ is the maximal cutoff radius for surgeries in
$[t_0/2, t_0]$. (If there are no surgeries on $[t_0/2, t_0]$ then we put
$h_{\max}(t_0) = 0$.)
Assume
\begin{enumerate}
\item The ball $B(x_0, t_0, r_0)$ has sectional curvatures at least $- r_0^{-2}$
at each point.
\item The volume of any subball $B(x, t_0, r) \subset B(x_0, t_0, r_0)$ with
any radius $r > 0$ is at least $(1-\epsilon)$ times the volume of the
Euclidean ball of the same radius.
\end{enumerate}

Then there is an unscathed time-dependent family
of metric balls
$\bigcup_{t \in [t_0-\tau^\prime r_0^2, t_0]} B(x_0, t, r_0/2)$
on which the supremum of $R$ is less than $K r_0^{-2}$.
\end{lemma}

\begin{proof}
\sketch Use the lower curvature and volume hypotheses with point selection, distance distortion, and local curvature estimates; a contradiction argument gives the stated upper bound.
\uses{lem:kl-curvature-bounds-from-priori-assumptions-estimates-curvature,lem:kl-earlier-scalar-curvature-bounds-smaller-balls-from-lower-curvature-bound-estimates-scalar-curvature,lem:kl-earlier-scalar-curvature-bounds-smaller-balls-from-lower-curvature-bound-estimates-scalar-curvature-4,lem:kl-length-distortion-estimates-distance-distortion-lemma,lem:kl-locating-small-balls-whose-subballs-almost-euclidean-volume-existence-sectional-curvature,prop:kl-proof-double-sided-curvature-bound-thick-part-modulo-two-propositions-existence-canonical-neighborhoods}
\end{proof}

\begin{lemma}[scalar curvature lemma]
  \label{lem:kl-earlier-scalar-curvature-bounds-smaller-balls-from-lower-curvature-bound-scalar-curvature-lemma}
  \source{kleiner-lott}
  \dcref{Sublemma 7}
  \group{Kleiner--Lott Long-Time Curvature}

$r_0^\al > r(t_0^\al)$ for all but a finite number of $\al$.
\end{lemma}

\begin{proof}
\sketch Use the lower curvature and volume hypotheses with point selection, distance distortion, and local curvature estimates; a contradiction argument gives the stated upper bound.
\uses{lem:kl-curvature-bounds-from-priori-assumptions-estimates-curvature,lem:kl-earlier-scalar-curvature-bounds-smaller-balls-from-lower-curvature-bound-scalar-curvature-lemma}
\proves{lem:kl-earlier-scalar-curvature-bounds-smaller-balls-from-lower-curvature-bound-existence-scalar-curvature-6}
\end{proof}

\begin{lemma}[estimates for scalar curvature]
  \label{lem:kl-earlier-scalar-curvature-bounds-smaller-balls-from-lower-curvature-bound-estimates-scalar-curvature-8}
  \source{kleiner-lott}
  \uses{lem:kl-earlier-scalar-curvature-bounds-smaller-balls-from-lower-curvature-bound-estimates-scalar-curvature}
  \dcref{Sublemma 8}
  \group{Kleiner--Lott Long-Time Curvature}

$\widehat{\tau}$ is bounded below by the parameter
$\tau_0$ of Lemma \ref{I.11.6(a)}, where we take $w = 1-\epsilon$ in
Lemma \ref{I.11.6(a)}.
\end{lemma}

\begin{proof}
\sketch Use the lower curvature and volume hypotheses with point selection, distance distortion, and local curvature estimates; a contradiction argument gives the stated upper bound.
\uses{lem:kl-earlier-scalar-curvature-bounds-smaller-balls-from-lower-curvature-bound-estimates-scalar-curvature,lem:kl-length-distortion-estimates-distance-distortion-lemma,lem:kl-locating-small-balls-whose-subballs-almost-euclidean-volume-existence-sectional-curvature,prop:kl-proof-double-sided-curvature-bound-thick-part-modulo-two-propositions-existence-canonical-neighborhoods}
\proves{lem:kl-earlier-scalar-curvature-bounds-smaller-balls-from-lower-curvature-bound-existence-scalar-curvature-6}
\proves{prop:kl-earlier-scalar-curvature-bounds-smaller-balls-from-lower-curvature-bound-existence-scalar-curvature}
\end{proof}

\begin{remark}[existence for scalar curvature]
\label{rem:kl-earlier-scalar-curvature-bounds-smaller-balls-from-lower-curvature-bound-existence-scalar-curvature}
\source{kleiner-lott}
\uses{prop:kl-proof-double-sided-curvature-bound-thick-part-modulo-two-propositions-existence-thick-thin-geometry}
\dcref{Remark 9}
\group{Kleiner--Lott Long-Time Curvature}


\end{remark}

\begin{remark}[structural remark on scalar curvature]
\label{rem:kl-earlier-scalar-curvature-bounds-smaller-balls-from-lower-curvature-bound-structural-remark-scalar-curvature}
\source{kleiner-lott}
\uses{cor:kl-double-sided-curvature-bound-thick-part-existence-ricci-flow-surgery}
\dcref{Remark 86.9}
\group{Kleiner--Lott Long-Time Curvature}


\label{hmax}
\end{remark}

\section{Long-time geometry}
\dcref{II.7--II.8}
\dcref{II.7.1} \label{II.7.1}

\begin{lemma}[monotonicity for noncollapsing]
  \label{lem:kl-noncollapsed-pointed-limits-hyperbolic-monotonicity-noncollapsing}
  \source{kleiner-lott}
  \dcref{Lemma 1}
  \group{Kleiner--Lott Long-Time Topology}

$V(t) \: \left( t + \frac14 \right)^{- \: \frac32}$ is nonincreasing in $t$.
\end{lemma}

\begin{proof}
\sketch Use the reduced-length and reduced-volume identities, the maximum principle, and the cited surgery estimates; the resulting barriers give the stated bound.
\uses{lem:kl-establishing-noncollapsing-presence-surgery-vanishing-noncollapsing}

\begin{equation} \label{dVdt}
\frac{dV}{dt} \: = \: - \: \int_M R \: dV \: \le \: - \: R_{min} \: V.
\end{equation}
\end{proof}

\begin{definition}[noncollapsing definition]
  \label{def:kl-noncollapsed-pointed-limits-hyperbolic-noncollapsing-definition}
  \source{kleiner-lott}
  \dcref{Definition 87.3}
  \group{Kleiner--Lott Long-Time Topology}
 \label{VR}
Put $\overline{V} \: = \: \lim_{t \rightarrow \infty} V(t) \: \left( t + \frac14 \right)^{- \: \frac32}$
and $\hat{R}(t) \: = \: R_{min}(t) \: V(t)^{\frac23}$.
\end{definition}

\begin{lemma}[noncollapsing lemma]
  \label{lem:kl-noncollapsed-pointed-limits-hyperbolic-noncollapsing-lemma}
  \source{kleiner-lott}
  \dcref{Lemma 87.4}
  \group{Kleiner--Lott Long-Time Topology}
 \label{Rhatev}
On any time interval which is free of singular times,
and on which   $R_{min}(t) \: \le \: 0$ for all $t$
(which we are assuming), we have
\begin{equation} \label{Rhat}
\frac{d\hat{R}}{dt} \: \ge \: \frac23 \: \hat{R} \: V^{-1} \:
\int_M (R_{min} - R) \: dV.
\end{equation}
\end{lemma}

\begin{proof}
\sketch Use the reduced-length and reduced-volume identities, the maximum principle, and the cited surgery estimates; the resulting barriers give the stated bound.
\end{proof}

\begin{corollary}[monotonicity for noncollapsing]
  \label{cor:kl-noncollapsed-pointed-limits-hyperbolic-monotonicity-noncollapsing}
  \source{kleiner-lott}
  \dcref{Corollary 87.7}
  \group{Kleiner--Lott Long-Time Topology}
 \label{Rmono}
If $R_{min}(t) \: \le \: 0$ for all $t$
(which we are assuming) then $\hat{R}(t)$ is nondecreasing.
\end{corollary}

\begin{proof}
\sketch Use the reduced-length and reduced-volume identities, the maximum principle, and the cited surgery estimates; the resulting barriers give the stated bound.
\uses{lem:kl-noncollapsed-pointed-limits-hyperbolic-noncollapsing-lemma}
\end{proof}

Put $\overline{R} \: = \: \lim_{t \rightarrow \infty} \hat{R}(t)$.

\begin{lemma}[noncollapsing lemma]
  \label{lem:kl-noncollapsed-pointed-limits-hyperbolic-noncollapsing-lemma-5}
  \source{kleiner-lott}
  \dcref{Lemma 87.8}
  \group{Kleiner--Lott Long-Time Topology}
 \label{RVlemma}
If $\overline{V} \: > \: 0$ then
$\overline{R} \: \overline{V}^{-2/3} \: = \: - \: \frac32$.
\end{lemma}

\begin{proof}
\sketch Use the reduced-length and reduced-volume identities, the maximum principle, and the cited surgery estimates; the resulting barriers give the stated bound.
\uses{lem:kl-establishing-noncollapsing-presence-surgery-vanishing-noncollapsing}
\end{proof}

\begin{proposition}[convergence for noncollapsing]
  \label{prop:kl-noncollapsed-pointed-limits-hyperbolic-convergence-noncollapsing}
  \source{kleiner-lott}
  \dcref{Proposition 87.11}
  \group{Kleiner--Lott Long-Time Topology}
 \label{propII.7.1}
Given the flow $\M$,
suppose that we have a sequence of parabolic neighborhoods
$P(x^\al, t^\al, r \sqrt{t^\al}, - r^2 t^\al)$, for
$t^\al \rightarrow \infty$ and some fixed $r \in (0,1)$, such that
the scalings of the parabolic neighborhoods with factor
$t^\al$ smoothly converge to some limit
solution $(\M_\infty, (\overline{x}, 1), g_\infty(\cdot))$
defined in a parabolic neighborhood
$P(\overline{x}, 1, r, -r^2)$. Then $g_\infty(t)$ has
constant sectional curvature $- \: \frac{1}{4t}$.
\end{proposition}

\begin{proof}
\sketch Use the reduced-length and reduced-volume identities, the maximum principle, and the cited surgery estimates; the resulting barriers give the stated bound.
\uses{lem:kl-noncollapsed-pointed-limits-hyperbolic-noncollapsing-lemma}

\begin{equation} \label{poscont}
\int_{B(x^\alpha, {s} t^\alpha, r \sqrt{t^\alpha})}
(R_{min}({s} t^\alpha) \: - \: R(x, {s} t^\alpha)) \:
dV \: < \: \frac{c}{2}
\: \sqrt{t^\alpha}.
\end{equation}

\begin{align} \label{Ralign}
\overline{R} \: - \: \hat{R}(0) \: & \ge \:
\frac23 \: \int_0^\infty \hat{R}(t) \: V(t)^{-1} \:
\int_M (R_{min}(t) - R(x,t)) \: dV(x) \: dt \\
& \ge \:
\frac23 \: \sum_{\al} t^\al
\int_{s_0- \mu}^{s_0 + \mu} \hat{R}(s t^\al) \: V(s t^\al)^{-1} \:
\int_M (R_{min}(s t^\al) - R(x,s t^\al)) \: dV(x) \: ds \notag \\
& \ge \:
\frac23 \: \sum_{\al} t^\al
\int_{s_0- \mu}^{s_0 + \mu} \hat{R}(s t^\al) \: V(s t^\al)^{-1} \:
\int_{B(x^\alpha, {s} t^\alpha, r \sqrt{t^\alpha})}
(R_{min}(s t^\al) - R(x,s t^\al)) \: dV(x) \: ds. \notag
\end{align}
\end{proof}

\subsection{Noncollapsed regions with a lower curvature bound are almost hyperbolic on a large scale}
\dcref{II.7.2}

\begin{lemma}[noncollapsing lemma]
  \label{lem:kl-noncollapsed-regions-lower-curvature-bound-almost-hyperbolic-large-scale-noncollapsing-lemma}
  \source{kleiner-lott}
  \dcref{Lemma II.7.2}
  \group{Kleiner--Lott Long-Time Topology}
 \label{lemII.7.2} (cf. Lemma II.7.2) \\
(a) Given $w, r, \xi > 0$ one can find $T = T(w,r,\xi) < \infty$ such that
if the ball $B(x_0, t_0, r \sqrt{t_0}) \subset \M_{t_0}^+$ at some time $t_0 \ge T$ has
volume at least $wr^3 t_0^{\frac32}$ and sectional curvatures at least
$- \: r^{-2} t_0^{-1}$ then the curvature at $(x_0, t_0)$
satisfies
\begin{equation} \label{closeness}
|2tR_{ij}(x_0, t_0) + g_{ij}|^2 \: = \:
(2tR_{ij}(x_0, t_0) + g_{ij})(2tR^{ij}(x_0, t_0) + g^{ij}) \: < \: \xi^2.
\end{equation}
(b) Given in addition $A < \infty$ and allowing $T$ to depend on $A$, we
can ensure (\ref{closeness}) for all points in
$B(x_0, t_0, Ar \sqrt{t_0})$. \\
(c) The same is true for $P(x_0, t_0, Ar \sqrt{t_0}, A r^2 t_0)$.

Note that the time $T$ will depend on the initial metric.
\end{lemma}

\begin{proof}
\sketch Use the reduced-length and reduced-volume identities, the maximum principle, and the cited surgery estimates; the resulting barriers give the stated bound.
\uses{cor:kl-double-sided-curvature-bound-thick-part-existence-ricci-flow-surgery,prop:kl-noncollapsed-pointed-limits-hyperbolic-convergence-noncollapsing,prop:kl-proof-double-sided-curvature-bound-thick-part-modulo-two-propositions-existence-canonical-neighborhoods}
\end{proof}

\subsection{Thick-thin decomposition}
\dcref{II.7.3}

\begin{definition}[curvature scale]
  \label{def:kl-thick-thin-decomposition-curvature-scale}
  \source{kleiner-lott}
  \dcref{Definition 1}
  \group{Kleiner--Lott Long-Time Topology}

For $x \in \M_t^+$,
let $\rho(x,t)$ be the unique number $\rho\in (0,\infty)$ such that
$\inf_{B^+(x,t, \rho)} \Rm \: = \: - \: \rho^{-2}$, if such a $\rho$ exists,
and put $\rho(x,t)=\infty$ otherwise.
\end{definition}

\begin{lemma}[thick--thin geometry lemma]
  \label{lem:kl-thick-thin-decomposition-thick-thin-geometry-lemma}
  \source{kleiner-lott}
  \dcref{Lemma 89.2}
  \group{Kleiner--Lott Long-Time Topology}
 \label{bigenough}
For any $w > 0$ we can find $\overline{\rho} = \overline{\rho}(w) > 0$
and $\overline{T}=\overline{T}(w)$
such that if $t\geq \overline{T}$ and $\rho(x,t) < \overline{\rho} \sqrt{t}$ then
\begin{equation}
\vol(B(x,t,\rho(x,t))) \: < \: w \rho^3(x,t).
\end{equation}
\end{lemma}

\begin{equation} \label{Rblowup}
\lim_{\al \rightarrow \infty}
R(x^{\al,\prime}, t^\al) \: \rho^2(x^\al, t^\al) \: = \: \infty.
\end{equation}

\begin{remark}[estimates for canonical neighborhoods]
\label{rem:kl-thick-thin-decomposition-estimates-canonical-neighborhoods}
\source{kleiner-lott}
\dcref{Remark 3}
\group{Kleiner--Lott Long-Time Topology}


The canonical-neighborhood and Bishop--Gromov argument gives the same contradiction after choosing the universal $\epsilon$ sufficiently small.
\end{remark}

\begin{definition}[thick--thin decomposition]
  \label{def:kl-thick-thin-decomposition-thick-thin-decomposition}
  \source{kleiner-lott}
  \dcref{Definition 4}
  \group{Kleiner--Lott Long-Time Topology}

The {\em $w$-thin part} $M^-(w, t) \subset \M_t^+$ is the set of points
$x \in M$ so that either $\rho(x,t)=\infty$ or
\begin{equation}
\vol(B(x, t, \rho(x, t))) \: < \: w \:
(\rho(x, t))^3.
\end{equation}
The {\em $w$-thick part} is $M^+(w, t) = \M_t^+  - M^-(w, t)$.
\end{definition}

\begin{lemma}[existence for thick--thin geometry]
  \label{lem:kl-thick-thin-decomposition-existence-thick-thin-geometry}
  \source{kleiner-lott}
  \uses{lem:kl-thick-thin-decomposition-thick-thin-geometry-lemma}
  \dcref{Lemma 89.10}
  \group{Kleiner--Lott Long-Time Topology}
 \label{applies}
Given $w > 0$, there are
$w^\prime = w^\prime(w) > 0$ and $T^\prime = T^\prime(w) < \infty$ so that
taking $r = \overline{\rho}(w)$
(with reference to Lemma \ref{bigenough}),
if $x_0 \in M^+(w,t)$ and $t_0 \ge T^\prime$ then
$B(x_0, t_0, r\sqrt{t_0})$ has volume at least $w^\prime r^3 t_0^{\frac32}$ and
sectional curvature at least $- \: r^{-2} \: t_0^{-1}$.
\end{lemma}

\begin{proof}
\sketch Apply the thick--thin decomposition and pointed smooth compactness. Hyperbolic rigidity and the approximation estimates give the required stabilization or geometric conclusion.
\uses{lem:kl-thick-thin-decomposition-thick-thin-geometry-lemma}
\end{proof}

\subsection{Hyperbolic rigidity and stabilization of the thick part}
\label{hyprig}

\begin{proposition}[monotonicity for thick--thin geometry]
  \label{prop:kl-hyperbolic-rigidity-stabilization-thick-part-monotonicity-thick-thin-geometry}
  \source{kleiner-lott}
  \dcref{Proposition 90.1}
  \group{Kleiner--Lott Long-Time Topology}

\label{slowisotopy}
There exist a number $T_0<\infty$, a nonincreasing function
$\alpha:[T_0,\infty)\ra (0,\infty)$ with
$\lim_{t \rightarrow \infty} \alpha(t) \: = \: 0$,
a (possibly empty) collection
$\{(H_1,x_1),\ldots,(H_k,x_k)\}$
of complete connected pointed finite-volume hyperbolic $3$-manifolds
and a family of smooth maps
\begin{equation}
f(t):B_t= \bigcup_{i=1}^k \;B \left( x_i, \frac{1}{\alpha(t)} \right)\lra \M_t,
\end{equation}
 defined for $t\in [T_0,\infty)$, such that

1.  $f(t)$ is close to an isometry:
\begin{equation}
\|t^{-1}f(t)^*g_{\M_t}-g_{B_t}\|_{C^{[\frac{1}{\alpha(t)}]}}<\alpha(t),
\end{equation}

2. $f(t)$ defines a smooth family of maps which changes slowly
with time:
\begin{equation}
|\dot{f}(p,t)|<\alpha(t)t^{-\frac12}
\end{equation}
for all $p\in B_t$,
where $\dot{f}$ refers to the time derivative (as defined
with admissible curves),

and

3. $f(t)$ parametrizes more and more of the thick part:
 $M^+(\alpha(t),t)\subset \im(f(t))$ for all $t\geq T_0$.
\end{proposition}

\begin{remark}[characterization for thick--thin geometry]
\label{rem:kl-hyperbolic-rigidity-stabilization-thick-part-characterization-thick-thin-geometry}
\source{kleiner-lott}
\uses{prop:kl-hyperbolic-rigidity-stabilization-thick-part-monotonicity-thick-thin-geometry}
\dcref{Remark 2}
\group{Kleiner--Lott Long-Time Topology}



For sufficiently small fixed $w$, the pointed-limit hyperbolic manifolds are independent of $w$ and their thick parts approximate the thick part of the flow.
\end{remark}

\begin{definition}[pointed smooth closeness]
  \label{def:kl-hyperbolic-rigidity-stabilization-thick-part-pointed-smooth-closeness}
  \source{kleiner-lott}
  \dcref{Equation 90.8}
  \group{Kleiner--Lott Long-Time Topology}

If $(X,x)$ and $(Y,y)$ are pointed smooth Riemannian manifolds
and $\eps>0$ then $(X,x)$ is $\eps$-close to $(Y,y)$
if there is a pointed map $f:(X,x)\ra (Y,y)$ such that
\begin{equation}
f\restr_{\ol{B(x,\eps^{-1})}}:\ol{B(x,\eps^{-1})}\ra Y
\end{equation}
 is a diffeomorphism onto its image
and
\begin{equation}
\label{fapprox}
\|f^*g_Y-g_X\|_{C^{\eps^{-1}}}<\eps,
\end{equation}
where the norm is taken on $\ol{B(x,\eps^{-1})}$.
Note that nothing is required of $f$ on the complement of $\ol{B(x,\eps^{-1})}$.
Such a map $f$ is called an {\em $\eps$-approximation}.
\end{definition}

\begin{lemma}[existence for thick--thin geometry]
  \label{lem:kl-hyperbolic-rigidity-stabilization-thick-part-existence-thick-thin-geometry}
  \source{kleiner-lott}
  \dcref{Lemma 90.11}
  \group{Kleiner--Lott Long-Time Topology}

\label{prasad}
Let $(X,x)$ be a pointed complete connected finite-volume hyperbolic
$3$-manifold.  Then for each $\zeta>0$ there exists $\xi>0$
such that if $X'$ is a complete finite-volume hyperbolic
manifold with at least as many cusps as $X$,
and $f:(X,x)\ra X'$ is a $\xi$-approximation, then there is an
isometry $\hat f: (X,x)\ra X'$ which
is $\zeta$-close to $f$.
\end{lemma}

\begin{proof}
\sketch Apply the thick--thin decomposition and pointed smooth compactness. Hyperbolic rigidity and the approximation estimates give the required stabilization or geometric conclusion.
The retained proof uses \cite{Mostow,Prasad}.
\end{proof}

\begin{definition}[hyperbolic pointed-limit space]
  \label{def:kl-hyperbolic-rigidity-stabilization-thick-part-hyperbolic-pointed-limit-space}
  \source{kleiner-lott}
  \dcref{Definition 5}
  \group{Kleiner--Lott Long-Time Topology}

Given $w >0$, let $\Lambda_w$ be the space
of complete pointed finite-volume hyperbolic
$3$-manifolds that arise as pointed limits of
sequences $\{(\M^+_{t_i}, (x_i, t_i), t_i^{-1} g(t_i))\}_{i=1}^\infty$
with $\lim_{i \rightarrow \infty} t_i = \infty$ and the basepoint
$(x_i,t_i)$ satisfying
$(x_i,t_i) \in M^+(w, t_i) \subset \M^+_{t_i}$
for all $i$.
\end{definition}

\begin{lemma}[existence for cap geometry]
  \label{lem:kl-hyperbolic-rigidity-stabilization-thick-part-existence-cap-geometry}
  \source{kleiner-lott}
  \dcref{Lemma 90.14}
  \group{Kleiner--Lott Long-Time Topology}

\label{II.7.2again}
Given $w>0$, there is a decreasing function
$\beta:[0,\infty)\ra (0,\infty]$
with $\lim_{s\ra\infty}\beta(s)=0$ such that
if
$(x,t)\in M^+(w,t)\subset\M_t^+$,
and  $Z_t$ denotes
the forward time slice
$\M_t^+$ rescaled by $t^{-1}$, then

1.  Some $(X,x)\in \La_w$
is $\beta(t)$-close to $(Z_t,(x,t))$.

2.
$B(x,t,\beta(t)^{-1}\sqrt{t})\subset \M_t^+$
is unscathed
on the interval $[t,2t]$ and if $\ga:[t,2t]\ra \M$ is a static
curve starting at $(x,t)$, $\bar t\in [t,2t]$, then the map
\begin{equation}
B(x,t,\beta(t)^{-1}\sqrt{t})\ra
P(x,t,\beta(t)^{-1}\sqrt{t},t)\cap\M_{\bar t}
\end{equation}
defined by following static curves induces a map
\begin{equation}
i_{t,\bar t}:(Z_t,(x,t))\supset (B(x,t,\beta(t)^{-1}),(x,t))
\ra (Z_{\bar t},\ga(\bar t))
\end{equation}
satisfying
\begin{equation}
\|\left(i_{t,\bar t}\right)^*g_{Z_{\bar
t}}-g_{Z_t}\|_{C^{\beta(t)^{-1}}}
<\beta(t).
\end{equation}
\end{lemma}

\begin{proof}
\sketch Apply the thick--thin decomposition and pointed smooth compactness. Hyperbolic rigidity and the approximation estimates give the required stabilization or geometric conclusion.
\uses{lem:kl-noncollapsed-regions-lower-curvature-bound-almost-hyperbolic-large-scale-noncollapsing-lemma}
\end{proof}

\subsection{Incompressibility of cuspidal tori} \label{inctori}
\dcref{II.7}

Let $T_a<\infty$ be large enough  that  $B_{T_a}$  (defined as in Proposition \ref{slowisotopy})
contains $Y$.

\begin{equation}
Z =  \partial V_{j},\quad\hat Z_t =  f(t)(Z)\\
\quad \hat Y_t =  f(t)(Y),\quad
\hat W_t =  \M_t^+ -  \Int(\hat Y_t)
\end{equation}

\begin{proposition}[incompressibility of cuspidal tori proposition]
  \label{prop:kl-incompressibility-cuspidal-tori-incompressibility-cuspidal-tori-proposition}
  \source{kleiner-lott}
  \dcref{Proposition 91.2}
  \group{Kleiner--Lott Long-Time Topology}

\label{propincompressible}
The homomorphism
\begin{equation} \label{injeqn}
\pi_1(f(t)):\pi_1(Z,\star)\ra \pi_1(\M_t^+,f(t)(\star))
\end{equation}
is a monomorphism for all $t\geq T_a$.
\end{proposition}

\begin{lemma}[incompressibility of cuspidal tori lemma]
  \label{lem:kl-incompressibility-cuspidal-tori-incompressibility-cuspidal-tori-lemma}
  \source{kleiner-lott}
  \dcref{Lemma 91.4}
  \group{Kleiner--Lott Long-Time Topology}
 \label{stablelemma}
The kernels of the homomorphisms
\begin{equation}
\label{eqnpi1inclusion}
\pi_1(f(t)):\pi_1(Z,\star)\ra \pi_1(\M_t^+,f(t)(\star)),\quad
\pi_1(f(t)):\pi_1(Z,\star)\ra \pi_1(\hat W_t,f(t)(\star))
\end{equation}
are independent of $t$, for all $t\geq T_a$.
\end{lemma}

\begin{equation}
(N^2_t,\D N^2_t)\subset (\hat W_t,\hat Z_t).
\end{equation}

\begin{lemma}[existence for Ricci flow with surgery]
  \label{lem:kl-incompressibility-cuspidal-tori-existence-ricci-flow-surgery}
  \source{kleiner-lott}
  \dcref{Lemma 91.10}
  \group{Kleiner--Lott Long-Time Topology}

\label{lemavoids}
Let $\de(t)$ be the surgery parameter from
Section \ref{II.4.4}. There is a $T = T(a) < \infty$ so that whenever $t \ge T$,
no point in any area-minimizing compressing disk  $N_t\subset \hat W_t$ is in the
center of a $10 \delta(t)$-neck.
\end{lemma}

\begin{proof}
\sketch Apply the preceding estimates and compactness argument to the stated hypotheses; the limiting or inductive construction has the required properties.
The retained proof uses \cite{schoen}, \cite[Theorem 2]{schoenyau}.
\end{proof}

\begin{lemma}[existence for kappa solutions]
  \label{lem:kl-incompressibility-cuspidal-tori-existence-kappa-solutions}
  \source{kleiner-lott}
  \dcref{Lemma 91.11}
  \group{Kleiner--Lott Long-Time Topology}
 \label{Hamlemma}
Given $D > 0$, there is a number $a_0 > 0$
so that for any $a \in (0,a_0)$, if $\diam(Z) = a$ and $t$ is
sufficiently large then
$\int_{\D N_t}\kappa_{\D N_t}ds \: \le \: \frac{D}{2}$ and
$\length(\D N_t) \: \le \: \frac{D}{2} \: \sqrt{t}$,
where $\kappa_{\D N_t}$ is the geodesic curvature of
$\D N_t \subset N_t$.
\end{lemma}

\begin{proof}
\sketch Argue by contradiction and rescale at the first failure. Hamilton--Ivey pinching, noncollapsing, and compactness yield a $\kappa$-solution or standard-solution limit, contradicting the failure.
\uses{prop:kl-hyperbolic-rigidity-stabilization-thick-part-monotonicity-thick-thin-geometry}
The retained proof uses \cite[Sections 11 and 12]{Hamiltonn}, \cite[Section 12]{Hamiltonn}.
\end{proof}

\begin{lemma}[existence for diameter]
  \label{lem:kl-incompressibility-cuspidal-tori-existence-diameter}
  \source{kleiner-lott}
  \dcref{Lemma 91.12}
  \group{Kleiner--Lott Long-Time Topology}
 \label{lemsupport}
For every $D>0$ there is a number $a_0>0$ with the following property.
Given $a \in (0, a_0)$, suppose that we take $Z$ to be the torus
cross-section
in $V_j$
of diameter $a$. Then
there is a number $T^\prime_a < \infty$ so that as long as
$t_0\geq T^\prime_a$, there is a smooth function $\bar A$ defined
on a neighborhood of $t_0$ such that $\bar  A(t_0)=A(t_0)$, $\bar A\geq A$
everywhere, and
\begin{equation} \label{dereqn}
\bar A'(t_0)< \frac{3}{4}\left(\frac{1}{t_0+\frac14}\right) A(t_0)-2\pi+D.
\end{equation}
\end{lemma}

\begin{proof}
\sketch Combine the canonical-neighborhood, noncollapsing, and distance/volume comparison estimates, treating surgery and nonsurgery cases separately.
\uses{lem:kl-establishing-noncollapsing-presence-surgery-vanishing-noncollapsing,lem:kl-incompressibility-cuspidal-tori-existence-kappa-solutions,lem:kl-incompressibility-cuspidal-tori-existence-ricci-flow-surgery,prop:kl-hyperbolic-rigidity-stabilization-thick-part-monotonicity-thick-thin-geometry}

\begin{equation}
\label{eqninitial}
\bar A'(t_0)=\int_{\D S}\langle X,\nu_{\D S}\rangle\:ds
+\int_S\frac{d}{dt} \: \Big|_{t=t_0} \dvol_S,
\end{equation}

\begin{equation}
\label{eqnbdyterm}
\left|\int_{\D S_{t_0}}\langle X,\nu_{\D S_{t_0}}\rangle\:ds\right| \:
\leq \: \alpha(t_0) \: t_0^{-\frac12} \: \length(\D N_{t_0}).
\end{equation}

\begin{equation} \label{evolve}
\int_S\:\frac{d}{dt} \: \Big|_{t=t_0} \dvol_S \: \le \:
\int_S\:\frac34\left(\frac{1}{t_0+\frac14}\right) \: \dvol_S+\int_{\D S}\kappa_{\D S}ds-2\pi.
\end{equation}
\end{proof}

\subsection{The thin part is a graph manifold}
\dcref{II.7.4} \label{II.7.4}

\begin{theorem}[existence for thick--thin geometry]
  \label{thm:kl-thin-part-graph-manifold-existence-thick-thin-geometry}
  \source{kleiner-lott}
  \dcref{Theorem II.7.4}
  \group{Kleiner--Lott Long-Time Topology}
 (cf. Theorem II.7.4) \label{thmII.7.4}
Suppose that $(M^\al, g^\al)$ is a sequence of compact oriented Riemannian
$3$-manifolds, closed or with convex boundary, and $w^\al \rightarrow 0$.
Assume that \\
(1) for each point $x \in M^\al$ there exists a radius $\rho = \rho^\al(x)$
not exceeding the diameter of $M^\al$
such that
the ball $B(x, \rho)$ in the metric $g^\al$ has volume at most
$w^\al \rho^3$ and sectional curvatures at least $- \: \rho^{-2}$.\\
(2) each component of the boundary of $M^\al$ has diameter at most
$w^\al$, and has a (topologically trivial) collar of length one, where
the sectional curvatures are between $- \: \frac14 \: - \: \epsilon$ and
$- \: \frac14 \: - \: \epsilon$.

Then for large $\alpha$, $M^\al$ is diffeomorphic to a graph manifold.
\end{theorem}

\begin{remark}[estimates for thick--thin geometry]
\label{rem:kl-thin-part-graph-manifold-estimates-thick-thin-geometry}
\source{kleiner-lott}
\uses{thm:kl-thin-part-graph-manifold-existence-thick-thin-geometry}
\dcref{Remark 2}
\group{Kleiner--Lott Long-Time Topology}



The graph-manifold theorem applies under the stated collar and derivative hypotheses; the derivative bound is needed only through the finite range $m\leq[\epsilon^{-1}]$.
\end{remark}

\begin{theorem}[existence for thick--thin geometry]
  \label{thm:kl-thin-part-graph-manifold-existence-thick-thin-geometry-3}
  \source{kleiner-lott}
  \dcref{Theorem II.7.4}
  \group{Kleiner--Lott Long-Time Topology}
 (cf. Theorem II.7.4) \label{thmII.7.4second}
Suppose that $(M^\al, g^\al)$ is a sequence of compact oriented Riemannian
$3$-manifolds, closed or with convex boundary, and $w^\al \rightarrow 0$.
Assume that \\
(1) for each point $x \in M^\al$ there exists a radius $\rho = \rho^\al(x)$
such that
the ball $B(x, \rho)$ in the metric $g^\al$ has volume at most
$w^\al \rho^3$ and sectional curvatures at least $- \: \rho^{-2}$.\\
(2) each component of the boundary of $M^\al$ has diameter at most
$w^\al$, and has a (topologically trivial) collar of length one, where
the sectional curvatures are between $- \: \frac14 \: - \: \epsilon$ and
$- \: \frac14 \: - \: \epsilon$. \\
(3) for every $w^\prime > 0$ and $m  \in \{0, 1, \ldots, [\epsilon^{-1}]\}$,
there exist $\overline{r} = \overline{r}(w^\prime) > 0$
and $K_m = K_m(w^\prime) <\infty$ such that for sufficiently
large $\alpha$, if $r \in (0, \overline{r}]$
and a ball $B(x,r)$ in the metric $g^\al$
has volume at least $w^\prime r^3$ and sectional curvatures at least
$- \: r^{-2}$ then $|\nabla^m \Rm|(x) \: \le \: K_m \: r^{-m-2}$.

Then for large $\alpha$, $M^\al$ is diffeomorphic to a graph manifold.
\end{theorem}

\begin{remark}[estimates for noncollapsing]
\label{rem:kl-thin-part-graph-manifold-estimates-noncollapsing}
\source{kleiner-lott}
\dcref{Remark 4}
\group{Kleiner--Lott Long-Time Topology}

\end{remark}

\begin{remark}[structural remark on thick--thin geometry]
\label{rem:kl-thin-part-graph-manifold-structural-remark-thick-thin-geometry}
\source{kleiner-lott}
\uses{thm:kl-thin-part-graph-manifold-existence-thick-thin-geometry-3}
\dcref{Remark 5}
\group{Kleiner--Lott Long-Time Topology}



Condition (3) supplies the derivative bounds used in the graph-manifold conclusion.
\end{remark}

\begin{remark}[positivity for thick--thin geometry]
\label{rem:kl-thin-part-graph-manifold-positivity-thick-thin-geometry}
\source{kleiner-lott}
\uses{thm:kl-thin-part-graph-manifold-existence-thick-thin-geometry, thm:kl-thin-part-graph-manifold-existence-thick-thin-geometry-3}
\dcref{Remark 6}
\group{Kleiner--Lott Long-Time Topology}



The radius bound in the first theorem distinguishes it from the second; without that bound the limit may be only a nonnegatively curved Alexandrov space.
\end{remark}

\begin{remark}[existence for thick--thin geometry]
\label{rem:kl-thin-part-graph-manifold-existence-thick-thin-geometry}
\source{kleiner-lott}
\uses{thm:kl-no-local-collapsing-theorem-ii-existence-kappa-solutions, thm:kl-thin-part-graph-manifold-existence-thick-thin-geometry, thm:kl-thin-part-graph-manifold-existence-thick-thin-geometry-3}
\dcref{Remark 7}
\group{Kleiner--Lott Long-Time Topology}



Homogeneous collapsed solutions satisfy the graph-manifold hypotheses, while finite-time noncollapsing is not contradicted by their long-time collapse.
\end{remark}

\begin{equation}
M_{thin}(t) \: = \: \M_t^+ - f(t) (\widehat{H}_1(t) \cup \ldots \widehat{H}_k(t)).
\end{equation}

\begin{proposition}[thick--thin geometry proposition]
  \label{prop:kl-thin-part-graph-manifold-thick-thin-geometry-proposition}
  \source{kleiner-lott}
  \dcref{Proposition 92.12}
  \group{Kleiner--Lott Long-Time Topology}
 \label{thingraph}
For large $t$, $M_{thin}(t)$ is a graph manifold.
\end{proposition}

\begin{proof}
\sketch Apply the thick--thin decomposition and pointed smooth compactness. Hyperbolic rigidity and the approximation estimates give the required stabilization or geometric conclusion.
\uses{cor:kl-double-sided-curvature-bound-thick-part-existence-ricci-flow-surgery,lem:kl-curvature-bounds-from-priori-assumptions-canonical-neighborhoods-lemma,lem:kl-curvature-bounds-from-priori-assumptions-estimates-curvature,lem:kl-locating-small-balls-whose-subballs-almost-euclidean-volume-existence-sectional-curvature,lem:kl-thin-part-graph-manifold-thick-thin-geometry-lemma,rem:kl-earlier-scalar-curvature-bounds-smaller-balls-from-lower-curvature-bound-structural-remark-scalar-curvature,thm:kl-thin-part-graph-manifold-existence-thick-thin-geometry,thm:kl-thin-part-graph-manifold-existence-thick-thin-geometry-3}
The retained proof uses \cite{Hamiltonnnnn}.
\end{proof}

\begin{lemma}[thick--thin geometry lemma]
  \label{lem:kl-thin-part-graph-manifold-thick-thin-geometry-lemma}
  \source{kleiner-lott}
  \uses{thm:kl-thin-part-graph-manifold-existence-thick-thin-geometry-3}
  \dcref{Lemma 92.13}
  \group{Kleiner--Lott Long-Time Topology}
 \label{condition}
Condition (3) in Theorem \ref{thmII.7.4second} holds for the $M^\al$'s.
\end{lemma}

\begin{proof}
\sketch Apply the thick--thin decomposition and pointed smooth compactness. Hyperbolic rigidity and the approximation estimates give the required stabilization or geometric conclusion.
\uses{cor:kl-double-sided-curvature-bound-thick-part-existence-ricci-flow-surgery,lem:kl-curvature-bounds-from-priori-assumptions-canonical-neighborhoods-lemma,lem:kl-curvature-bounds-from-priori-assumptions-estimates-curvature,lem:kl-locating-small-balls-whose-subballs-almost-euclidean-volume-existence-sectional-curvature,rem:kl-earlier-scalar-curvature-bounds-smaller-balls-from-lower-curvature-bound-structural-remark-scalar-curvature,thm:kl-thin-part-graph-manifold-existence-thick-thin-geometry-3}
\proves{prop:kl-thin-part-graph-manifold-thick-thin-geometry-proposition}
\end{proof}

\section{Alternative proof of cusp incompressibility}
\dcref{II.8} \label{II.8}

\begin{equation} \label{Rayleigh}
\lambda(g) \: = \: \inf_{\Phi \in C^\infty(M) \: : \:
\Phi \neq 0} \frac{\int_M \left( 4 |\nabla \Phi|^2 \: + \:
R \: \Phi^2 \right) \: dV}{\int_M \Phi^2 \: dV}.
\end{equation}

\begin{equation} \label{ll}
\frac{d}{dt} \lambda(t) \: \ge \: \frac23 \:
\lambda^2(t).
\end{equation}

For any metric $g$, there are inequalities

\begin{equation} \label{test}
\min R \: \le \: \lambda(g) \: \le \: \frac{\int_M R \: dV}{\vol(M, g)},
\end{equation}

\subsection{The approach using the $\sigma$-invariant}

\begin{equation} \label{yamabe}
\sigma(M) \: = \: \sup_{{\mathcal C}} \inf_{g \in {\mathcal C}}
\frac{\int_M R(g) \: \dvol(g)}{\vol(M, g)^{\frac{n-2}{n}}},
\end{equation}

\begin{equation} \label{testt}
\sigma(M) \: = \: \sup_{g} R_{min}(g) \: V(g)^{\frac2n}.
\end{equation}

\begin{proposition}[characterization for graph manifolds]
  \label{prop:kl-alternative-proof-cusp-incompressibility-characterization-graph-manifolds}
  \source{kleiner-lott}
  \dcref{Proposition 93.9}
  \group{Kleiner--Lott Long-Time Topology}
 \label{propYamabe}
Let $M$ be a closed connected oriented $3$-manifold. \\
(a) If $\sigma(M) > 0$ then $M$ is
diffeomorphic to a connected sum of a finite number of
$S^1 \times S^2$'s and metric quotients of the round $S^3$.
Conversely, each such manifold has $\sigma(M) > 0$. \\
(b) $M$ is a graph manifold if and only if $\sigma(M)
\ge 0$. \\
(c) If $\sigma(M) < 0$
 then
$\left( - \: \frac23 \: \sigma(M)
\right)^{\frac32}$ is the minimum
of the numbers
$V$ with the following property : $M$ can be decomposed as a connected sum
of a finite collection of $S^1 \times S^2$'s, metric quotients of the round
$S^3$ and some other components, the union of which is denoted by
$M^\prime$, and there exists a (possibly disconnected) complete
finite-volume
manifold $N$ with constant sectional curvature $- \: \frac14$ and
volume $V$ which can be embedded in $M^\prime$ so that the
complement $M^\prime - N$ (if nonempty) is a graph manifold.

Moreover, if
$\vol(N) \: = \: \left( - \: \frac23 \: \sigma(M)
\right)^{\frac32}$
then the
cusps of $N$ (if any) are incompressible in $M^\prime$.
\end{proposition}

\begin{proof}
\sketch Combine the Rayleigh-quotient or sigma-invariant inequalities with monotonicity under the flow and the long-time thick--thin decomposition; the extremal volume characterization follows.
\uses{lem:kl-double-sided-curvature-bound-thick-part-compactness-ricci-flow-surgery,lem:kl-double-sided-curvature-bound-thick-part-compactness-thick-thin-geometry}
The retained proof uses \cite[Pf. of Theorem 2.9]{Anderson}.
\end{proof}

\subsection{The approach using the $\overline{\lambda}$-invariant}

\begin{proposition}[characterization for graph manifolds]
  \label{prop:kl-alternative-proof-cusp-incompressibility-characterization-graph-manifolds-2}
  \source{kleiner-lott}
  \dcref{Proposition II.8.2}
  \group{Kleiner--Lott Long-Time Topology}
 (cf. II.8.2) \label{propII.8.2}
Let $M$ be a closed connected oriented $3$-manifold. \\
(a) If $M$ admits a metric $g$ with $\lambda(g) > 0$ then it is
diffeomorphic to a connected sum of a finite number of
$S^1 \times S^2$'s and metric quotients of the round $S^3$.
Conversely, each such manifold admits a metric $g$ with
$\lambda(g) > 0$. \\
(b) Suppose that $M$ does not admit any metric $g$ with
$\lambda(g) > 0$. Let $\overline{\lambda}$ denote the
supremum of $\lambda(g) V(g)^{\frac23}$ over all metrics $g$ on $M$.
Then $M$ is a graph manifold if and only if $\overline{\lambda}
= 0$. \\
(c) Suppose that $M$ does not admit any metric $g$ with
$\lambda(g) > 0$, and $\overline{\lambda} < 0$.
Then
$\left( - \: \frac23 \: \overline{\lambda}
\right)^{\frac32}$ is the minimum
of the numbers
$V$ with the following property : $M$ can be decomposed as a connected sum
of a finite collection of $S^1 \times S^2$'s, metric quotients of the round
$S^3$ and some other components, the union of which is denoted by
$M^\prime$, and there exists a (possibly disconnected) complete
manifold $N$ with constant sectional curvature $- \: \frac14$ and
volume $V$ which can be embedded in $M^\prime$ so that the
complement $M^\prime - N$ (if nonempty) is a graph manifold.

Moreover, if
$\vol(N) \: = \: \left( - \: \frac23 \: \overline{\lambda}
\right)^{\frac32}$
then the cusps of $N$ (if any) are incompressible in $M^\prime$.
\end{proposition}

\begin{proof}
\sketch Combine the Rayleigh-quotient or sigma-invariant inequalities with monotonicity under the flow and the long-time thick--thin decomposition; the extremal volume characterization follows.
\uses{lem:kl-alternative-proof-cusp-incompressibility-estimates-spectral-geometry,lem:kl-alternative-proof-cusp-incompressibility-positivity-entropy,lem:kl-double-sided-curvature-bound-thick-part-compactness-ricci-flow-surgery,lem:kl-establishing-noncollapsing-presence-surgery-vanishing-noncollapsing,prop:kl-statement-existence-theorem-ricci-flow-surgery-monotonicity-kappa-solutions}
The retained proof uses \cite[Pf. of Theorem 2.9]{Anderson}.
\end{proof}

\begin{lemma}[estimates for spectral geometry]
  \label{lem:kl-alternative-proof-cusp-incompressibility-estimates-spectral-geometry}
  \source{kleiner-lott}
  \dcref{Lemma 93.14}
  \group{Kleiner--Lott Long-Time Topology}
 \label{newprop1}
Given a closed Riemannian manifold $M$, let $X$ be a codimension-$0$
submanifold-with-boundary of $M$. Given $R \in C^\infty(M)$, let
$\lambda_M$ be the lowest eigenvalue of $-4\triangle \: + \: R$
on $M$, with corresponding eigenfunction $\psi$.
Let $\lambda_X$ be the lowest eigenvalue of the corresponding
operator on $X$, with Dirichlet boundary conditions, and similarly
for $\lambda_{M - \Int(X)}$. Then for all
$\eta \in C^\infty_c(\Int(X))$, we have
\begin{equation} \label{eig1}
\lambda_M \: \le \: \min(\lambda_X, \lambda_{M - \Int(X)})
\end{equation}
and
\begin{equation} \label{eig2}
\lambda_X \: \le \: \lambda_M \: + \: 4 \:
\frac{\int_M |\nabla \eta|^2 \: \psi^2 \: dV}{\int_M \eta^2 \:
\psi^2 \: dV}.
\end{equation}
\end{lemma}

\begin{proof}
\sketch Apply the preceding estimates and compactness argument to the stated hypotheses; the limiting or inductive construction has the required properties.
\proves{prop:kl-alternative-proof-cusp-incompressibility-characterization-graph-manifolds-2}
The retained proof uses \cite[Chapter XIII.15]{Reed-Simon}.
\end{proof}

\begin{lemma}[positivity for entropy]
  \label{lem:kl-alternative-proof-cusp-incompressibility-positivity-entropy}
  \source{kleiner-lott}
  \uses{lem:kl-alternative-proof-cusp-incompressibility-estimates-spectral-geometry}
  \dcref{Lemma 93.19}
  \group{Kleiner--Lott Long-Time Topology}
 \label{newprop2}
With the notation of Lemma \ref{newprop1},
given a nonnegative function $\phi \in C^\infty(M)$,
suppose that $f \in C^\infty(M)$ satisfies
\begin{equation}
4 |\nabla f|^2 \: \le \: R \: - \: \lambda_M \: - \: c
\end{equation}
on $\supp(\phi)$, for some $c > 0$. Then
\begin{equation} \label{eig3}
\|e^f \phi \psi\|_2 \: \le \: 4 \: c^{-1} \: \left( \|e^f \triangle \phi\|_\infty
\: + \: \|e^f \nabla \phi\|_\infty (\lambda_M \: - \: \min R)^{1/2} \right)
\: \|\psi\|_2.
\end{equation}
\end{lemma}

\begin{proof}
\sketch Use the reduced-length and reduced-volume identities, the maximum principle, and the cited surgery estimates; the resulting barriers give the stated bound.
\proves{prop:kl-alternative-proof-cusp-incompressibility-characterization-graph-manifolds-2}

\begin{equation} \label{part1}
\int_M e^f \psi\phi \: (H \: - \: 4 |\nabla f|^2 \: - \: \lambda_M) \: \phi
e^f \psi dV \: \ge \: c \int_M \phi^2 e^{2f} \psi^2 \: dV.
\end{equation}

\begin{equation} \label{Hf}
e^f H (e^{-f} \rho) \: = \: H \rho \: + \: 4 \: \nabla \cdot
((\nabla f) \rho) \: + \: 4 \: \langle \nabla f, \nabla \rho \rangle
\: - \: 4 \: |\nabla f|^2 \rho
\end{equation}

\begin{equation} \label{part2}
\int_M e^{2f} \phi \psi H(\phi \psi) \: dV \: = \:
\int_M e^f \phi \psi (H \: - \: 4 \:
|\nabla f|^2)  e^{f} \phi \psi \: dV.
\end{equation}

\begin{equation} \label{bound11}
\lambda_X \: \le \: \lambda^+ \: + \: \const \: h(T_0)^{-2} \:
\frac{\int_{N_h M_{cap}} \psi^2 \: dV}{1 - \int_{N_h M_{cap}} \psi^2 \: dV}.
\end{equation}

\begin{equation} \label{bound12}
\int_{N_h M_{cap}} \psi^2 \: dV \: \le \:
\left( \sup_{N_h M_{cap}} e^{-2f} \right) \:
\int_{N_h M_{cap}} \: e^{2f} \psi^2 \: dV \: \le \:
\const \: \sup_{N_h M_{cap}} e^{-2f}.
\end{equation}
\end{proof}

\begin{remark}[positivity for scalar curvature]
\label{rem:kl-alternative-proof-cusp-incompressibility-positivity-scalar-curvature}
\source{kleiner-lott}
\dcref{Remark 5}
\group{Kleiner--Lott Long-Time Topology}


If $M$ admits no metric of positive scalar curvature, then $\sigma(M)=\overline{\lambda}$.
\end{remark}
